\editionPage[\pdfbookmark[0]{Corrigés des défis}{bm-corriges-defis}]{defi-cor-1-1}{D01 · Corrigé}{Mélanger n'est pas moyenner}{\hyperref[defi-ex-1]{Énoncé p.~\pageref*{defi-ex-1}} · Démarche et vérifications}
\begin{blocEdition}{defi-cor-1-1}
\editionEtape{1. Additionner des quantités, pas des pH}
Les quantités d'ions oxonium valent $\num{1.0e-2}\times\num{0.1000}=\qty{1,0e-3}{mol}$ et
$\num{1.0e-4}\times\num{0.1000}=\qty{1,0e-5}{mol}$. Dans les 200,0 mL du mélange,
\[\hho=\frac{\qty{1,01e-3}{mol}}{\qty{0,2000}{L}}=\qty{5,05e-3}{mol.L^{-1}}\qquad\text{donc}\qquad\res{\pH\approx2{,}30}.\]
Le mélange est bien plus proche de la solution la plus acide que la moyenne 3,0 ne le laisserait penser :
la solution de pH 2,0 apporte cent fois plus d'ions oxonium que l'autre.
\editionEtape{2. Une moyenne de concentrations, pas de pH}
Avec $a=10^{-p_1}$ et $b=10^{-p_2}$, le mélange à volumes égaux a pour concentration réduite $h=\frac{a+b}{2}$.
L'inégalité admise donne $h\geq\sqrt{ab}=10^{-(p_1+p_2)/2}$. La fonction $x\mapsto-\log x$ étant
strictement décroissante,
\[p=-\log h\leq-\log\sqrt{ab}=\frac{p_1+p_2}{2}.\]
L'égalité a lieu si et seulement si $a=b$, c'est-à-dire $p_1=p_2$. Le pH d'un mélange est toujours
inférieur ou égal à la moyenne des pH.
\editionEtape{3. Ajouter un acide pour diluer}
Soit $V$ le volume ajouté, en litres. On veut
\[\frac{\num{1.0e-3}+\num{1.0e-4}\,V}{0{,}1000+V}=\num{1.0e-3}\iff\num{9.0e-4}=\num{9.0e-4}\,V
\iff\res{V=\qty{1,0}{L}}.\]
Avec de l'eau pure, $\frac{\num{1.0e-3}}{0{,}1000+V'}=\num{1.0e-3}$ donne $V'=\qty{0,90}{L}$.
Il faut davantage de solution de pH 4,0 que d'eau : elle dilue, mais elle apporte aussi ses propres ions oxonium.
\editionEtape{4. Tenir compte des ions de l'eau}
La relation simple donnerait $\pH=8$ : une solution d'acide serait basique, ce qui est absurde.
À cette dilution, les ions oxonium apportés par l'acide ne sont plus prépondérants devant ceux de l'eau.
On pose $x=c/\cst$ ; on a $[\ce{Cl-}]=c$ et $[\ce{HO-}]=\dfrac{K_\mathrm{e}\,{\cst}^2}{\hho}$, soit, en valeurs
réduites, $\dfrac{K_\mathrm{e}}{h}$. L'électroneutralité donne $h=x+\dfrac{K_\mathrm{e}}{h}$, soit
$h^2-xh-K_\mathrm{e}=0$. La racine positive est
\[h=\frac{x+\sqrt{x^2+4K_\mathrm{e}}}{2}=\num{1.05e-7}\qquad\text{donc}\qquad\res{\pH\approx6{,}98}.\]
\end{blocEdition}

\editionPage{defi-cor-1-2}{D01 · Corrigé (suite)}{Mélanger n'est pas moyenner}{\hyperref[defi-ex-1]{Énoncé p.~\pageref*{defi-ex-1}} · Démarche et vérifications}
\begin{blocEdition}{defi-cor-1-2}
\editionEtape{5. Domaine de validité de la relation simple}
L'écart vaut $\delta=-\log x-(-\log h)=\log\dfrac{h}{x}=\log\dfrac{1+\sqrt{1+4K_\mathrm{e}/x^2}}{2}$ ;
il décroît quand $x$ augmente. Alors
\[\delta<0{,}01\iff\sqrt{1+\frac{4K_\mathrm{e}}{x^2}}<2\times10^{0{,}01}-1\iff
x>\sqrt{\frac{4K_\mathrm{e}}{(2\times10^{0{,}01}-1)^2-1}}\approx\num{6.5e-7}.\]
\begin{codecote}[0.6]
\begin{pythoncol}
from math import sqrt, log10
KE = 1.0e-14
def ph_exact(c):
    h = (c + sqrt(c**2 + 4*KE)) / 2
    return -log10(h)
print(round(ph_exact(1e-8), 3))
\end{pythoncol}
\cote
\begin{consolecol}
6.978
\end{consolecol}
\end{codecote}
\label{figedition-defi-cor-1-2-1}\schema[corrige]{defi-01-courbe}
\editionRetenir{Le pH est une grandeur logarithmique : on ne moyenne pas des pH, on écrit la conservation
des quantités de matière. Un modèle a aussi un domaine de validité : la relation simple convient tant que
l'acide apporte beaucoup plus d'ions oxonium que l'eau, c'est-à-dire pour $c$ supérieure à environ \qty{6,5e-7}{mol.L^{-1}}.}
\end{blocEdition}

\editionPage{defi-cor-2-1}{D02 · Corrigé}{La méthode des ajouts dosés}{\hyperref[defi-ex-2]{Énoncé p.~\pageref*{defi-ex-2}} · Démarche et vérifications}
\begin{blocEdition}{defi-cor-2-1}
\editionEtape{1. L'effet de matrice}
Dans l'échantillon, d'autres espèces peuvent interagir avec X ou modifier son environnement : le coefficient $k$
n'est alors plus celui mesuré dans l'eau pure. La droite d'étalonnage « idéale » ne s'applique pas à l'échantillon ;
les ajouts dosés, réalisés dans l'échantillon lui-même, mesurent le $k$ de la matrice réelle.
\editionEtape{2. Concentration dans chaque fiole}
La quantité de X vaut $n=c_0V_0+c_\mathrm{s}V_\mathrm{s}$, répartie dans le volume $V_\mathrm{T}$ :
\[c=\frac{c_0V_0+c_\mathrm{s}V_\mathrm{s}}{V_\mathrm{T}}\qquad\text{donc}\qquad
A=k\,c=\frac{k}{V_\mathrm{T}}\left(c_0V_0+c_\mathrm{s}V_\mathrm{s}\right).\]
\editionEtape{3. Une droite et son intersection avec l'axe}
$A=a\,V_\mathrm{s}+b$ avec $a=\dfrac{k\,c_\mathrm{s}}{V_\mathrm{T}}$ et $b=\dfrac{k\,c_0V_0}{V_\mathrm{T}}$ : la
fonction est affine. Elle s'annule pour
\[V_\mathrm{s}^{*}=-\frac{b}{a}=-\frac{c_0V_0}{c_\mathrm{s}}\qquad\text{d'où}\qquad c_0=\frac{b}{a}\times\frac{c_\mathrm{s}}{V_0}.\]
Le coefficient $k$ s'élimine : il n'est pas nécessaire de le connaître.
\editionEtape{4. Exploiter les mesures}
\begin{codecote}[0.62]
\begin{pythoncol}
import numpy as np
Vs = np.array([0, 2.0, 4.0, 6.0, 8.0])
A = np.array([0.209, 0.311, 0.417, 0.520, 0.623])
a, b = np.polyfit(Vs, A, 1)
c0 = b / a * 1.00e-4 / 10.0
print(round(a, 5), round(b, 4), f"{c0:.2e}")
\end{pythoncol}
\cote
\begin{consolecol}
0.05185 0.2086 4.02e-05
\end{consolecol}
\end{codecote}
$\frac{b}{a}=\qty{4,02}{mL}$ et $\res{c_0=\qty{4,0e-5}{mol.L^{-1}}}$. Enfin, $k=\dfrac{a\,V_\mathrm{T}}{c_\mathrm{s}}
=\dfrac{\qty{0,05185}{mL^{-1}}\times\qty{50,0}{mL}}{\qty{1,00e-4}{mol.L^{-1}}}\approx\res{\qty{2,6e4}{L.mol^{-1}}}$.
\end{blocEdition}

\editionPage{defi-cor-2-2}{D02 · Corrigé (suite)}{La méthode des ajouts dosés}{\hyperref[defi-ex-2]{Énoncé p.~\pageref*{defi-ex-2}} · Démarche et vérifications}
\begin{blocEdition}{defi-cor-2-2}
\label{figedition-defi-cor-2-2-1}\schema[corrige]{defi-02-ajouts}
\editionEtape{5. Une absorbance parasite}
L'absorbance parasite s'ajoute à toutes les mesures : la pente $a$ ne change pas, mais l'ordonnée à l'origine
mesurée vaut $b'=b+A_\mathrm{p}$. La concentration calculée à la question 4 (avec $b'=0{,}2086$) est donc trop
grande ; la vraie concentration correspond à $b=0{,}2086-0{,}020=0{,}1886$ :
\[c_0=\frac{0{,}1886}{0{,}05185}\times\frac{\num{1.00e-4}}{10{,}0}=\res{\qty{3,6e-5}{mol.L^{-1}}}.\]
En ignorant l'absorbance parasite, on surestime $c_0$ d'environ \res{11~\%}. On l'évite en mesurant un « blanc »
contenant la matrice sans l'espèce X, ou l'absorbance de fond à une longueur d'onde où X n'absorbe pas, puis en
soustrayant cette valeur de toutes les mesures.
\editionRetenir{Une méthode d'analyse repose sur un modèle : il faut savoir ce qu'il suppose (ici, la même
valeur de $k$ dans toutes les fioles et l'absence d'absorbance parasite) et contrôler ces hypothèses.}
\end{blocEdition}

\editionPage{defi-cor-3-1}{D03 · Corrigé}{Identifier un acide inconnu}{\hyperref[defi-ex-3]{Énoncé p.~\pageref*{defi-ex-3}} · Démarche et vérifications}
\begin{blocEdition}{defi-cor-3-1}
\editionEtape{1. Une famille confirmée}
La bande très large entre 2\,500 et 3\,200 cm$^{-1}$ correspond à la liaison O--H d'un acide carboxylique, et la bande
forte vers 1\,700 cm$^{-1}$ à la liaison C=O : le solide est bien un acide carboxylique.
\editionEtape{2. L'hypothèse d'un monoacide}
$n(\ce{HO-})=0{,}100\times\num{25.4e-3}=\qty{2,54e-3}{mol}$. S'il n'y avait qu'un groupe carboxyle,
$n(\text{acide})=\qty{2,54e-3}{mol}$ et $M=\frac{0{,}150}{\num{2.54e-3}}=\qty{59,1}{g.mol^{-1}}$. L'analyse élémentaire donne alors :
\[x=\frac{0{,}407\times59{,}1}{12{,}0}=2{,}0\ ;\qquad y=\frac{0{,}051\times59{,}1}{1{,}0}=3{,}0\ ;\qquad z=\frac{0{,}542\times59{,}1}{16{,}0}=2{,}0.\]
La formule \ce{C2H3O2} comporte un nombre impair d'atomes d'hydrogène : chaque atome de carbone forme quatre liaisons,
chaque oxygène deux et chaque hydrogène une ; la somme $4x+2z+y$ vaut le double du nombre de liaisons, donc elle est
paire, et $y$ doit être pair. \res{\text{L'hypothèse d'un monoacide est à rejeter.}}
\editionEtape{3. Le nombre de groupes carboxyle}
Avec $p$ groupes, $n(\text{acide})=\frac{\num{2.54e-3}}{p}$ et $M=p\times\qty{59,1}{g.mol^{-1}}$ : la formule est
$(\ce{C2H3O2})_p$, soit $\ce{C}_{2p}\ce{H}_{3p}\ce{O}_{2p}$. Elle compte $\frac{4p+2-3p}{2}=\frac{p+2}{2}$ doubles liaisons,
qui doivent être au moins $p$ (une par groupe C=O) : $p+2\geq2p$, donc $p\leq2$. Comme $p=1$ est exclu, $\res{p=2}$ :
\[\res{\ce{C4H6O4},\quad M=\qty{118,1}{g.mol^{-1}}}.\]
On vérifie : $\frac{4\times12{,}01}{118{,}1}=40{,}7~\%$ de carbone et $n(\text{acide})=\frac{0{,}150}{118{,}1}=\qty{1,27e-3}{mol}$,
exactement la moitié de $n(\ce{HO-})$.
\end{blocEdition}

\editionPage{defi-cor-3-2}{D03 · Corrigé (suite)}{Identifier un acide inconnu}{\hyperref[defi-ex-3]{Énoncé p.~\pageref*{defi-ex-3}} · Démarche et vérifications}
\begin{blocEdition}{defi-cor-3-2}
\editionEtape{4. Choisir entre deux isomères}
Deux diacides du tableau ont pour formule brute \ce{C4H6O4} : l'acide méthylpropanedioïque et l'acide
butanedioïque. La température de fusion mesurée, 184 °C, désigne l'\res{\text{acide butanedioïque}}
\ce{HOOC-CH2-CH2-COOH}, aussi appelé acide succinique. L'analyse élémentaire ne donne que la formule brute :
elle ne distingue pas des isomères, d'où la nécessité d'une mesure physique supplémentaire.
\editionEtape{5. Un seul saut pour deux acidités}
\label{figedition-defi-cor-3-2-1}\schema[corrige]{defi-03-courbe}
Les deux groupes carboxyle cèdent leur ion \ce{H+} dans des conditions voisines : la seconde acidité commence à
réagir avant que la première soit entièrement titrée. Les deux réactions se recouvrent et le pH ne varie
brutalement qu'une fois, lorsque les deux acidités ont réagi :
\[\ce{H2A + 2 HO- -> A^2- + 2 H2O},\qquad n(\ce{HO-})_\mathrm{E}=2\,n(\text{acide}).\]
C'est pourquoi le titrage, interprété seul, conduisait à une masse molaire deux fois trop faible.
\editionRetenir{Un titrage compte les ions \ce{H+} cédés, pas les molécules : il faut connaître le nombre d'acidités
titrées. Croiser plusieurs méthodes (titrage, analyse élémentaire, spectroscopie, température de fusion) et vérifier la
cohérence chimique d'une formule évite les conclusions hâtives.}
\end{blocEdition}

\editionPage{defi-cor-4-1}{D04 · Corrigé}{Ordre 1 ou ordre 2 ?}{\hyperref[defi-ex-4]{Énoncé p.~\pageref*{defi-ex-4}} · Démarche et vérifications}
\begin{blocEdition}{defi-cor-4-1}
\editionEtape{1. Le test des demi-réactions successives}
La concentration vaut $c_0/2=\qty{5,00}{mmol.L^{-1}}$ vers $t=\qty{900}{s}$ (on mesure 5,03 mmol$\cdot$L$^{-1}$) :
$t_{1/2}\approx\qty{9e2}{s}$. Pour une loi d'ordre 1, la concentration serait à nouveau divisée par deux au bout de
900 s supplémentaires : on mesurerait $c_0/4=\qty{2,50}{mmol.L^{-1}}$ à $t=\qty{1800}{s}$. On mesure
\qty{3,36}{mmol.L^{-1}}, nettement plus : \res{\text{la loi de vitesse n'est pas d'ordre 1}}.
\editionEtape{2. La linéarisation de l'ordre 1 échoue}
\begin{center}\small\renewcommand{\arraystretch}{1.2}
\begin{tabular}{@{}l*{7}{c}@{}}\toprule
$t$ (s) & 0 & 120 & 300 & 600 & 900 & 1\,200 & 1\,800\\
$\ln(c/c^\circ)$ & $-4{,}605$ & $-4{,}730$ & $-4{,}890$ & $-5{,}113$ & $-5{,}292$ & $-5{,}447$ & $-5{,}696$\\\bottomrule
\end{tabular}\end{center}
Les coefficients directeurs successifs valent, en $10^{-4}$ s$^{-1}$ : $-10{,}4$ ; $-8{,}9$ ; $-7{,}4$ ; $-6{,}0$ ; $-5{,}1$ ;
$-4{,}1$. Ils ne sont pas égaux : les points ne sont pas alignés, $\ln(c/c^\circ)$ n'est pas une fonction affine du
temps. La concentration décroît moins vite qu'une exponentielle, ce qui confirme la question 1.
\editionEtape{3. Le modèle d'ordre 2}
\begin{codecote}[0.62]
\begin{pythoncol}
import numpy as np
t = np.array([0, 120, 300, 600, 900, 1200, 1800])
c = np.array([10.00, 8.83, 7.52, 6.02,
              5.03, 4.31, 3.36]) * 1e-3
k, b = np.polyfit(t, 1 / c, 1)
print(round(k, 4), round(b, 1))
\end{pythoncol}
\cote
\begin{consolecol}
0.1098 100.1
\end{consolecol}
\end{codecote}
Les valeurs de $1/c$ (100,0 ; 113,3 ; 133,0 ; 166,1 ; 198,8 ; 232,0 ; 297,6 L$\cdot$mol$^{-1}$) sont alignées
(graphique page suivante) : le modèle convient, avec $\res{k\approx\qty{0,11}{L.mol^{-1}.s^{-1}}}$. L'ordonnée à
l'origine, \qty{100}{L.mol^{-1}}, est bien égale à $1/c_0$.
\end{blocEdition}

\editionPage{defi-cor-4-2}{D04 · Corrigé (suite)}{Ordre 1 ou ordre 2 ?}{\hyperref[defi-ex-4]{Énoncé p.~\pageref*{defi-ex-4}} · Démarche et vérifications}
\begin{blocEdition}{defi-cor-4-2}
\label{figedition-defi-cor-4-2-1}\schema[corrige]{defi-04-cor}
\editionEtape{4. Un temps de demi-réaction qui dépend de $c_0$}
À la date $t_{1/2}$, $c=c_0/2$ : $\dfrac{2}{c_0}=\dfrac{1}{c_0}+k\,t_{1/2}$, d'où
\[t_{1/2}=\frac{1}{k\,c_0}=\frac{1}{0{,}11\times\num{1.00e-2}}\approx\res{\qty{9,1e2}{s}},\]
en accord avec l'estimation de la question 1. Pour $c_0=\qty{5,0e-3}{mol.L^{-1}}$, $t_{1/2}\approx\res{\qty{1,8e3}{s}}$ :
contrairement à l'ordre 1, le temps de demi-réaction double quand la concentration initiale est divisée par deux.
\editionEtape{5. Des demi-réactions de plus en plus longues}
Pour passer de $c_0/2$ à $c_0/4$, on applique le résultat précédent avec $c_0/2$ comme concentration initiale :
la durée vaut $\dfrac{1}{k\,c_0/2}=2\,t_{1/2}\approx\qty{1,8e3}{s}$. La concentration vaut donc $c_0/4$ à la date
\[t=t_{1/2}+2\,t_{1/2}=3\,t_{1/2}\approx\res{\qty{2,7e3}{s}}\qquad\text{(vérification : }\tfrac{1}{c}=100+0{,}11\times2\,727\approx\qty{400}{L.mol^{-1}}\text{)}.\]
Avec une loi d'ordre 1 de même temps de demi-réaction, $c_0/4$ serait atteinte dès $2\,t_{1/2}\approx\qty{1,8e3}{s}$.
L'ordre 2 ralentit davantage en fin de réaction : la vitesse $k\,c^2$ diminue plus vite que $c$ lorsque les réactifs
se raréfient.
\editionRetenir{Pour identifier une loi de vitesse, on teste une linéarisation : $\ln c$ en fonction de $t$ pour
l'ordre 1, $1/c$ pour l'ordre 2 (relation fournie). Les temps de demi-réaction successifs sont une signature
simple : constants pour l'ordre 1, ils doublent à chaque étape pour l'ordre 2.}
\end{blocEdition}

\editionPage{defi-cor-5-1}{D05 · Corrigé}{Une horloge à iode}{\hyperref[defi-ex-5]{Énoncé p.~\pageref*{defi-ex-5}} · Démarche et vérifications}
\begin{blocEdition}{defi-cor-5-1}
\editionEtape{1. Le diiode produit avant l'apparition de la couleur}
Jusqu'à $t_\mathrm{c}$, tout le diiode formé par (1) est consommé par (2), qui met en jeu deux ions thiosulfate par
molécule de diiode. À la date $t_\mathrm{c}$, les ions thiosulfate sont épuisés, donc
\[[\ce{I2}]_\text{produit}=\frac{[\ce{S2O3^2-}]_0}{2}=\res{\qty{2,5e-4}{mol.L^{-1}}}.\]
\editionEtape{2. Une vitesse pratiquement constante}
La réaction (1) consomme deux ions iodure par molécule de diiode formée, et la réaction (2) les restitue aussitôt :
$[\ce{I-}]$ \res{\text{reste égale à }[\ce{I-}]_0} jusqu'à $t_\mathrm{c}$. L'eau oxygénée, elle, n'est pas
régénérée : sa concentration diminue de \qty{2,5e-4}{mol.L^{-1}}, soit au plus $\frac{\num{2.5e-4}}{0{,}010}=2{,}5~\%$
(expérience 4). Les concentrations des réactifs de (1) étant pratiquement constantes, sa vitesse l'est aussi :
\[v=\frac{[\ce{I2}]_\text{produit}}{t_\mathrm{c}}=\frac{\qty{2,5e-4}{mol.L^{-1}}}{t_\mathrm{c}}.\]
\editionEtape{3. Tester la loi de vitesse}
Si $v=k\,[\ce{H2O2}]_0\,[\ce{I-}]_0$, alors $k=\dfrac{\num{2.5e-4}}{t_\mathrm{c}\,[\ce{H2O2}]_0\,[\ce{I-}]_0}$ :
\begin{center}\small\renewcommand{\arraystretch}{1.2}
\begin{tabular}{@{}l*{5}{c}@{}}\toprule
Expérience & 1 & 2 & 3 & 4 & 5\\\midrule
$k$ (L$\cdot$mol$^{-1}\cdot$s$^{-1}$) & 0,0120 & 0,0120 & 0,0116 & 0,0121 & 0,0119\\\bottomrule
\end{tabular}\end{center}
Les valeurs sont égales aux incertitudes près : doubler l'une ou l'autre des concentrations divise $t_\mathrm{c}$ par
deux. La loi est vérifiée, avec $\res{k\approx\qty{0,012}{L.mol^{-1}.s^{-1}}}$.
\label{figedition-defi-cor-5-1-1}\schema[corrige]{defi-05-cor}
\end{blocEdition}

\editionPage{defi-cor-5-2}{D05 · Corrigé (suite)}{Une horloge à iode}{\hyperref[defi-ex-5]{Énoncé p.~\pageref*{defi-ex-5}} · Démarche et vérifications}
\begin{blocEdition}{defi-cor-5-2}
\editionEtape{4. Un mécanisme cohérent}
En additionnant les actes a, b, c et d, les entités \ce{IO-}, \ce{HIO} et \ce{HO-} apparaissent une fois de chaque
côté ; il reste deux ions iodure, deux ions oxonium et quatre molécules d'eau :
\[\res{\ce{H2O2 + 2 I- + 2 H3O+ -> I2 + 4 H2O}}.\] Les intermédiaires réactionnels sont
\res{\ce{IO-}, \ce{HIO}\text{ et }\ce{HO-}}, formés puis consommés. La vitesse est celle de l'acte lent a, qui fait se
rencontrer une molécule de peroxyde d'hydrogène et un ion iodure : $v=k_\mathrm{a}\,[\ce{H2O2}]\,[\ce{I-}]$. C'est
\res{\text{la loi observée}} ; les ions oxonium, qui n'interviennent qu'après l'acte lent, n'y figurent pas.
\editionEtape{5. Prévoir, puis comprendre le rôle du thiosulfate}
\[t_\mathrm{c}=\frac{\num{2.5e-4}}{k\,[\ce{H2O2}]_0\,[\ce{I-}]_0}=\frac{\num{2.5e-4}}{0{,}012\times0{,}030\times0{,}030}\approx\res{\qty{23}{s}}.\]
Sans ions thiosulfate, le diiode s'accumulerait dès le début : la solution bleuirait progressivement, sans
changement brutal, et il n'y aurait plus d'horloge. Les ions thiosulfate fixent la quantité de diiode à produire
avant le virage et maintiennent constante la concentration en ions iodure.
\editionRetenir{Une horloge chimique mesure la durée nécessaire pour produire une quantité fixée de produit : si la
vitesse est constante, $1/t_\mathrm{c}$ lui est proportionnelle. Comparer des expériences qui ne diffèrent que par une
concentration permet de tester une loi de vitesse, puis de discuter un mécanisme.}
\end{blocEdition}

\editionPage{defi-cor-6-1}{D06 · Corrigé}{Le générateur de technétium}{\hyperref[defi-ex-6]{Énoncé p.~\pageref*{defi-ex-6}} · Démarche et vérifications}
\begin{blocEdition}{defi-cor-6-1}
\editionEtape{1. La désintégration du molybdène 99}
\[\res{\ce{^{99}_{42}Mo -> ^{99\text{m}}_{43}Tc + ^{0}_{-1}e}}\qquad\text{(nucléons : }99=99+0\text{ ; charge : }42=43+(-1)\text{)}.\]
\editionEtape{2. Un noyau qui se forme et se désintègre}
Pendant $\mathrm{d}t$, $\lambda_1N_1\,\mathrm{d}t$ noyaux de molybdène se désintègrent et donnent autant de noyaux de
technétium ; dans le même temps, $\lambda_2N_2\,\mathrm{d}t$ noyaux de technétium se désintègrent. Ainsi
$\mathrm{d}N_2=\lambda_1N_1\,\mathrm{d}t-\lambda_2N_2\,\mathrm{d}t$, soit $\res{\dfrac{\mathrm{d}N_2}{\mathrm{d}t}=\lambda_1N_1-\lambda_2N_2}$.
\editionEtape{3. Vérifier la solution}
On pose $K=\frac{\lambda_1N_{1,0}}{\lambda_2-\lambda_1}$. Alors
$\frac{\mathrm{d}N_2}{\mathrm{d}t}=K\left(-\lambda_1\mathrm{e}^{-\lambda_1t}+\lambda_2\mathrm{e}^{-\lambda_2t}\right)$ et
\[\lambda_1N_1-\lambda_2N_2=\lambda_1N_{1,0}\,\mathrm{e}^{-\lambda_1t}-\lambda_2K\left(\mathrm{e}^{-\lambda_1t}-\mathrm{e}^{-\lambda_2t}\right)
=K\left[(\lambda_2-\lambda_1)\,\mathrm{e}^{-\lambda_1t}-\lambda_2\mathrm{e}^{-\lambda_1t}+\lambda_2\mathrm{e}^{-\lambda_2t}\right],\]
qui est bien égal à $\frac{\mathrm{d}N_2}{\mathrm{d}t}$. Enfin, $N_2(0)=K(1-1)=0$ : \res{\text{la condition initiale est respectée}}.
\editionEtape{4. Le meilleur moment pour éluer}
$A_2=\lambda_2N_2$ est maximale quand $\frac{\mathrm{d}N_2}{\mathrm{d}t}=0$, c'est-à-dire quand $\lambda_1N_1=\lambda_2N_2$ :
à cette date, \res{A_1=A_2}. D'après l'expression précédente, $\lambda_2\mathrm{e}^{-\lambda_2t}=\lambda_1\mathrm{e}^{-\lambda_1t}$,
soit $\mathrm{e}^{(\lambda_2-\lambda_1)t}=\frac{\lambda_2}{\lambda_1}$, d'où l'expression de $t_\mathrm{m}$. Avec
$\lambda_1=\frac{\ln2}{66{,}0}=\qty{0,01050}{h^{-1}}$ et $\lambda_2=\frac{\ln2}{6{,}01}=\qty{0,1153}{h^{-1}}$ :
\[t_\mathrm{m}=\frac{\ln(0{,}1153/0{,}01050)}{0{,}1153-0{,}01050}\approx\res{\qty{23}{h}}.\]
Au bout d'environ un jour, la quantité de technétium récupérable est maximale : une élution quotidienne, le matin,
fournit le technétium de la journée.
\end{blocEdition}

\editionPage{defi-cor-6-2}{D06 · Corrigé (suite)}{Le générateur de technétium}{\hyperref[defi-ex-6]{Énoncé p.~\pageref*{defi-ex-6}} · Démarche et vérifications}
\begin{blocEdition}{defi-cor-6-2}
\label{figedition-defi-cor-6-2-1}\schema[corrige]{defi-06-cor}
\editionEtape{5. L'équilibre transitoire}
\[\frac{A_2}{A_1}=\frac{\lambda_2N_2}{\lambda_1N_1}=\frac{\lambda_2}{\lambda_2-\lambda_1}\times\frac{\mathrm{e}^{-\lambda_1t}-\mathrm{e}^{-\lambda_2t}}{\mathrm{e}^{-\lambda_1t}}
=\frac{\lambda_2}{\lambda_2-\lambda_1}\left(1-\mathrm{e}^{-(\lambda_2-\lambda_1)t}\right).\]
Quand $t$ est grand devant $t_\mathrm{m}$, l'exponentielle tend vers 0 et
\[\frac{A_2}{A_1}\longrightarrow\frac{\lambda_2}{\lambda_2-\lambda_1}=\frac{0{,}1153}{0{,}1048}\approx\res{1{,}10}.\]
Le technétium, dont la demi-vie est onze fois plus courte que celle du molybdène, finit par décroître au rythme de
son père : les deux activités diminuent ensemble avec la demi-vie du molybdène, dans un rapport constant. Le
générateur reste ainsi utilisable environ une semaine, son activité étant divisée par deux tous les 66~h.
\editionRetenir{Lorsqu'un noyau radioactif est lui-même produit par un autre noyau radioactif, son évolution est
régie par une équation différentielle qui couple les deux populations. Si le père vit beaucoup plus longtemps que
le fils, les deux activités finissent par décroître ensemble : c'est le principe des générateurs utilisés en médecine.}
\end{blocEdition}

\editionPage{defi-cor-7-1}{D07 · Corrigé}{Le critère d'évolution vu par les vitesses}{\hyperref[defi-ex-7]{Énoncé p.~\pageref*{defi-ex-7}} · Démarche et vérifications}
\begin{blocEdition}{defi-cor-7-1}
\editionEtape{1. Deux actes qui se compensent}
L'espèce C est formée par l'acte direct, à la vitesse $v_1$, et consommée par l'acte inverse, à la vitesse $v_{-1}$ :
\[\res{\frac{\mathrm{d}[\mathrm{C}]}{\mathrm{d}t}=v_1-v_{-1}=k_1[\mathrm{A}][\mathrm{B}]-k_{-1}[\mathrm{C}][\mathrm{D}]}.\]
\editionEtape{2. La constante d'équilibre}
À l'équilibre dynamique, les deux vitesses sont égales : $k_1[\mathrm{A}]_\text{éq}[\mathrm{B}]_\text{éq}=k_{-1}[\mathrm{C}]_\text{éq}[\mathrm{D}]_\text{éq}$,
d'où
\[K=Q_{\mathrm{r},\text{éq}}=\frac{[\mathrm{C}]_\text{éq}[\mathrm{D}]_\text{éq}}{[\mathrm{A}]_\text{éq}[\mathrm{B}]_\text{éq}}=\frac{k_1}{k_{-1}}=\frac{0{,}020}{0{,}0050}=\res{4{,}0}.\]
\editionEtape{3. Le signe de la vitesse}
En factorisant $v_1$ :
\[\frac{\mathrm{d}[\mathrm{C}]}{\mathrm{d}t}=k_1[\mathrm{A}][\mathrm{B}]\left(1-\frac{k_{-1}}{k_1}\times\frac{[\mathrm{C}][\mathrm{D}]}{[\mathrm{A}][\mathrm{B}]}\right)
=k_1[\mathrm{A}][\mathrm{B}]\left(1-\frac{Q_\mathrm{r}}{K}\right).\]
Le facteur $k_1[\mathrm{A}][\mathrm{B}]$ est positif, donc :
\begin{itemize}[leftmargin=1.4em,itemsep=0.1em]
\item si $Q_\mathrm{r}<K$, la concentration de C augmente : le système évolue dans le sens direct ;
\item si $Q_\mathrm{r}>K$, elle diminue : le système évolue dans le sens inverse ;
\item si $Q_\mathrm{r}=K$, elle ne varie plus : le système est à l'équilibre.
\end{itemize}
C'est \res{\text{le critère d'évolution spontanée}}, qui découle ici du modèle cinétique.
\editionEtape{4. Une évolution dans le sens inverse}
$Q_{\mathrm{r},\mathrm{i}}=\dfrac{1{,}50\times1{,}50}{0{,}50\times0{,}50}=9{,}0>K$ : \res{\text{le système évolue dans le sens inverse}}.
Si $y$ est la quantité de C consommée, en mol, l'équilibre vérifie
\[\frac{(1{,}50-y)^2}{(0{,}50+y)^2}=4{,}0,\quad\text{soit}\quad\frac{1{,}50-y}{0{,}50+y}=2{,}0\quad\text{et}\quad
y=\frac16\approx\qty{0,17}{mol},\]
ce qui donne la composition du système à l'équilibre :
\[n(\mathrm{A})=n(\mathrm{B})=\res{\qty{0,667}{mol}},\qquad n(\mathrm{C})=n(\mathrm{D})=\res{\qty{1,333}{mol}}.\]
Le programme affiche 0,666~mol, 1,334~mol et $Q_\mathrm{r}=4{,}01$ : au bout de 150~min, l'équilibre est presque atteint.
\end{blocEdition}

\editionPage{defi-cor-7-2}{D07 · Corrigé (suite)}{Le critère d'évolution vu par les vitesses}{\hyperref[defi-ex-7]{Énoncé p.~\pageref*{defi-ex-7}} · Démarche et vérifications}
\begin{blocEdition}{defi-cor-7-2}
\label{figedition-defi-cor-7-2-1}\schema[corrige]{defi-07-cor}
\editionEtape{5. Le rôle d'un catalyseur}
Si $k_1$ et $k_{-1}$ sont multipliés par le même facteur, leur quotient ne change pas : \res{K\text{ est inchangée}}, ainsi que
la composition du système à l'équilibre. Les deux vitesses sont plus grandes : l'équilibre est atteint
\res{\text{plus rapidement}}. Un catalyseur accélère l'évolution vers l'état d'équilibre sans le déplacer.
\editionRetenir{Le modèle de l'équilibre dynamique relie la cinétique à l'équilibre : la constante d'équilibre est le
quotient des constantes de vitesse des deux actes opposés, et le signe de la vitesse globale dépend de la position
de $Q_\mathrm{r}$ par rapport à $K$. Un catalyseur modifie la durée de l'évolution, jamais l'état final.}
\end{blocEdition}

\editionPage{defi-cor-8-1}{D08 · Corrigé}{Quand une pile est-elle usée ?}{\hyperref[defi-ex-8]{Énoncé p.~\pageref*{defi-ex-8}} · Démarche et vérifications}
\begin{blocEdition}{defi-cor-8-1}
\editionEtape{1. La tension initiale}
\[\res{\ce{Cu^2+(aq) + Zn(s) -> Cu(s) + Zn^2+(aq)}},\qquad Q_\mathrm{r}=\frac{[\ce{Zn^2+}]}{[\ce{Cu^2+}]}.\]
Initialement, $Q_\mathrm{r}=\frac{0{,}10}{0{,}10}=1$, donc $\log Q_\mathrm{r}=0$ et $\res{U=\qty{1,10}{V}}$.
\editionEtape{2. Tension nulle et constante d'équilibre}
La tension à vide s'annule lorsque $0{,}0296\log Q_\mathrm{r}=1{,}10$, soit $\log Q_\mathrm{r}=\frac{1{,}10}{0{,}0296}\approx37{,}2$. La
pile ne peut plus débiter lorsque le système a atteint l'équilibre, c'est-à-dire lorsque $Q_\mathrm{r}=K$ : on en déduit
\[K\approx10^{37{,}2}\approx\res{\num{1.5e37}}.\]
C'est le même ordre de grandeur que \num{1.9e37} : l'écart provient de l'arrondi de la tension initiale, une valeur de
1,103~V redonnant exactement \num{1.9e37}.
\editionEtape{3. Composition et capacité}
Lorsque la charge $q$ a circulé, $n(\mathrm{e}^-)=q/F$ et l'avancement vaut $x=\frac{q}{2F}$, car deux électrons sont
échangés par ion \ce{Cu^2+} réduit :
\[[\ce{Cu^2+}]=c-\frac{q}{2FV},\qquad[\ce{Zn^2+}]=c+\frac{q}{2FV}.\]
Les ions \ce{Cu^2+} sont épuisés lorsque $q=Q_{\max}=2FcV=2\times\num{96500}\times0{,}10\times0{,}100=\res{\qty{1930}{C}}$,
soit \qty{0,54}{A.h}.
\editionEtape{4. Une tension qui varie peu}
À mi-décharge, $Q_\mathrm{r}=\frac{0{,}15}{0{,}05}=3$ et $U=\qty{1,086}{V}$. Lorsque 99,99~\% de la capacité a été utilisée,
la tension vaut encore \qty{0,973}{V}. Elle ne perd que quelques centièmes de volt jusqu'aux derniers instants :
\res{\text{la tension à vide ne permet pas de repérer une pile presque usée}}.
\end{blocEdition}

\editionPage{defi-cor-8-2}{D08 · Corrigé (suite)}{Quand une pile est-elle usée ?}{\hyperref[defi-ex-8]{Énoncé p.~\pageref*{defi-ex-8}} · Démarche et vérifications}
\begin{blocEdition}{defi-cor-8-2}
\label{figedition-defi-cor-8-2-1}\schema[corrige]{defi-08-cor}
\editionEtape{5. Ce qui fixe la durée de vie}
Si $U=0$, alors $Q_\mathrm{r}=K$ et
\[[\ce{Cu^2+}]=\frac{[\ce{Zn^2+}]}{K}\approx\frac{0{,}20}{\num{1.5e37}}\approx\res{\qty{1e-38}{mol.L^{-1}}}.\]
Dans 100~mL, cela représente $N=\num{1e-38}\times0{,}100\times\num{6.02e23}\approx\num{8e-16}$ ion, c'est-à-dire beaucoup moins
qu'un ion : l'équilibre ne serait atteint qu'après la réduction du tout dernier ion \ce{Cu^2+}. En pratique, la pile
cesse de débiter lorsque les ions \ce{Cu^2+} sont épuisés : \res{\text{le réactif limitant fixe sa capacité et sa durée
de vie}}. La baisse de tension observée en fin de vie d'une pile réelle vient de sa résistance interne et de la
lenteur avec laquelle les derniers ions atteignent l'électrode.
\editionRetenir{Pour une pile dont la constante d'équilibre est très grande, l'équilibre n'est jamais atteint en
pratique : la pile s'arrête quand son réactif limitant est épuisé. C'est pourquoi sa capacité se calcule à partir de sa
constitution initiale, avec $Q=n\,x_{\max}\,F$.}
\end{blocEdition}

\editionPage{defi-cor-9-1}{D09 · Corrigé}{Le pH d'un acide très dilué}{\hyperref[defi-ex-9]{Énoncé p.~\pageref*{defi-ex-9}} · Démarche et vérifications}
\begin{blocEdition}{defi-cor-9-1}
\editionEtape{1. Un acide ne rend jamais une solution basique}
L'eau pure contient déjà des ions \ce{H3O+} et \ce{HO-} en quantités égales ($\pH=7{,}0$ à 25~°C). Ajouter un acide ne
peut qu'\textbf{augmenter} la concentration en ions oxonium : le pH ne peut que diminuer. La valeur « $\pH=8$ » vient
d'une relation utilisée hors de son domaine de validité : elle néglige les ions \ce{H3O+} apportés par l'eau, qui sont
ici plus nombreux que ceux apportés par l'acide.
\editionEtape{2. L'électroneutralité de la solution}
La solution est électriquement neutre : $[\ce{H3O+}]=[\ce{Cl-}]+[\ce{HO-}]$. Avec $[\ce{Cl-}]=c$ et, d'après le produit
ionique, $[\ce{HO-}]=K_\mathrm{e}/h$ :
\[h=c+\frac{K_\mathrm{e}}{h}\qquad\text{soit}\qquad\res{h^2-c\,h-K_\mathrm{e}=0}.\]
\editionEtape{3. La racine positive}
Le discriminant vaut $c^2+4K_\mathrm{e}>0$ ; la seule racine positive est
\[h=\frac{c+\sqrt{c^2+4K_\mathrm{e}}}{2},\]
ce qui est exactement la ligne \code{h = (c + sqrt(c**2 + 4*Ke)) / 2} du programme. Pour $c=\qty{1,0e-8}{mol.L^{-1}}$ :
\[h=\frac{\num{1.0e-8}+\sqrt{\num{1.0e-16}+\num{4.0e-14}}}{2}=\qty{1,05e-7}{mol.L^{-1}},\qquad\pH=\res{6{,}98}.\]
La solution est très légèrement acide, comme attendu.
\editionEtape{4. Deux régimes}
Pour $c=\qty{1,0e-2}{mol.L^{-1}}$ et $c=\qty{1,0e-6}{mol.L^{-1}}$, le programme retrouve $\pH=-\log(c/\cst)$ (2,0 et
6,0) : tant que $c\gtrsim\qty{1e-6}{mol.L^{-1}}$, les ions apportés par l'acide dominent et \res{\text{la relation
usuelle est fiable}}. En dessous, l'autoprotolyse domine : le pH se rapproche de 7,0 (6,98 à $\num{1e-8}$, 7,0 à
$\num{1e-10}$) ; \res{\text{le pH tend vers 7,0}} quand $c$ tend vers 0, la solution tend vers l'eau pure.
\end{blocEdition}

\editionPage{defi-cor-9-2}{D09 · Corrigé (suite)}{Le pH d'un acide très dilué}{\hyperref[defi-ex-9]{Énoncé p.~\pageref*{defi-ex-9}} · Démarche et vérifications}
\begin{blocEdition}{defi-cor-9-2}
\label{figedition-defi-cor-9-2-1}\schema[corrige]{defi-09-cor}
\editionEtape{5. Le cas d'une base forte très diluée}
Pour une base forte de concentration apportée $c$, l'électroneutralité s'écrit $[\ce{HO-}]=c+K_\mathrm{e}/[\ce{HO-}]$ :
c'est la même équation, vérifiée cette fois par $\omega=[\ce{HO-}]$. On en déduit
\[\omega=\frac{c+\sqrt{c^2+4K_\mathrm{e}}}{2}\qquad\text{puis}\qquad\pH=\mathrm{p}K_\mathrm{e}-\mathrm{pOH}=14{,}0+\log\frac{\omega}{\cst}.\]
Le pH suit $14{,}0+\log(c/\cst)$ tant que $c\gtrsim\qty{1e-6}{mol.L^{-1}}$, puis \res{\text{tend vers 7,0 par valeurs
supérieures}} : la courbe est la symétrique de celle de l'acide par rapport à la droite horizontale $\pH=7{,}0$.
\editionRetenir{Une relation comme $\pH=-\log(c/\cst)$ n'est valable que dans son domaine : elle suppose que l'acide
apporté impose la concentration en ions oxonium. Dès que l'apport devient comparable à celui de l'autoprotolyse de
l'eau, c'est l'équation complète, issue de l'électroneutralité et du produit ionique, qui s'applique : le pH de toute
solution très diluée tend vers celui de l'eau pure.}
\end{blocEdition}

\editionPage{defi-cor-10-1}{D10 · Corrigé}{Pourquoi au moins 1,23 volt ?}{\hyperref[defi-ex-10]{Énoncé p.~\pageref*{defi-ex-10}} · Démarche et vérifications}
\begin{blocEdition}{defi-cor-10-1}
\editionEtape{1. Un bras de fer avec la pile}
Dès que du dihydrogène et du dioxygène couvrent les électrodes, la cellule est aussi une pile de tension à vide
$U_0=\qty{1,23}{V}$, qui tend à faire circuler le courant dans le sens de la transformation spontanée
\ce{2 H2 + O2 -> 2 H2O}. Pour imposer le sens inverse, le générateur doit l'emporter sur cette pile :
\res{\text{il faut }U>U_0}. En dessous, le courant d'électrolyse ne peut pas s'établir durablement.
\editionEtape{2. L'énergie électrique minimale par mole}
Produire une mole de \ce{H2} à la cathode (\ce{2 H2O + 2 e- -> H2 + 2 HO-}) exige le passage de deux moles
d'électrons, soit la charge $Q=2F$. L'énergie électrique reçue vaut $W=Q\,U$ ; elle est minimale pour $U=U_0$ :
\[W_0=2F\,U_0=2\times\num{96500}\times1{,}23=\qty{2,37e5}{J}\approx\res{\qty{237}{kJ}}.\]
\editionEtape{3. Par kilogramme}
Un kilogramme de dihydrogène représente $\dfrac{1000}{2{,}0}=500$ mol, d'où
\[W=500\times\num{2.374e5}=\qty{1,19e8}{J}=\frac{\num{1.19e8}}{\num{3.6e6}}\approx\res{\qty{33,0}{kW.h}}.\]
Aucun électrolyseur, aussi perfectionné soit-il, ne descendra sous cette valeur.
\editionEtape{4. Le prix de la tension}
La deuxième colonne donne l'énergie électrique réellement consommée par mole, $2F\,U$ : elle croît avec la tension
(251~kJ à 1,3~V, 463~kJ à 2,4~V). Le quotient $U_0/U$ est le \res{\text{rendement énergétique}} de la cellule : la part
de l'énergie électrique réellement stockée sous forme chimique. Il \res{\text{décroît}} quand $U$ augmente : à 2,0~V,
tension des cellules industrielles, il ne vaut plus que 0,61.
\end{blocEdition}

\editionPage{defi-cor-10-2}{D10 · Corrigé (suite)}{Pourquoi au moins 1,23 volt ?}{\hyperref[defi-ex-10]{Énoncé p.~\pageref*{defi-ex-10}} · Démarche et vérifications}
\begin{blocEdition}{defi-cor-10-2}
\label{figedition-defi-cor-10-2-1}\schema[corrige]{defi-10-cor}
\editionEtape{5. Une électrolyse qui aspirerait de la chaleur}
La combustion d'une mole de dihydrogène restitue 242~kJ, alors que l'électrolyse idéale n'en reçoit que 237 du
générateur : il manque 5~kJ. Une cellule fonctionnant exactement à 1,23~V devrait donc \res{\text{absorber de la
chaleur}} du milieu extérieur pour boucler son bilan d'énergie : elle se refroidirait. La cellule n'échange aucune
chaleur lorsque l'énergie électrique couvre exactement les 242~kJ, c'est-à-dire pour la tension « thermoneutre »
\[U=\frac{\num{2.42e5}}{2F}=\res{\qty{1,25}{V}}.\]
Au-delà, l'excédent est dissipé : à 2,0~V, chaque mole produite dégage $386-242=\qty{144}{kJ}$ de chaleur, qu'il faut
évacuer des stacks industriels.
\editionRetenir{Une électrolyse n'est possible que si le générateur impose une tension supérieure à la tension à vide
de la pile correspondante : l'énergie électrique minimale pour former les produits vaut $Q\,U_0$. Toute tension
supplémentaire est dissipée en chaleur, et le rendement d'une cellule réelle se lit simplement dans le rapport
$U_0/U$.}
\end{blocEdition}

\editionPage{defi-cor-11-1}{D11 · Corrigé}{Combien d'isomères ?}{\hyperref[defi-ex-11]{Énoncé p.~\pageref*{defi-ex-11}} · Démarche et vérifications}
\begin{blocEdition}{defi-cor-11-1}
\editionEtape{1. Autant de liaisons que de carbones, moins une}
Des isomères de constitution ont la même formule brute mais des enchaînements d'atomes différents. Dans un squelette
non cyclique, chaque nouvel atome de carbone se raccorde à la structure existante par exactement une liaison C--C :
en partant d'un premier carbone (0 liaison), les $n-1$ suivants apportent chacun une liaison, d'où
\res{n-1\text{ liaisons C--C}}, quelle que soit la forme du squelette.
\editionEtape{2. Les trois squelettes de C5H12}
On classe selon la longueur de la plus longue chaîne.
\begin{itemize}[leftmargin=1.4em,itemsep=0.1em]
\item Chaîne de 5 : le \res{\text{pentane}}.
\item Chaîne de 4 : un méthyle à placer. En position 1, il allongerait la chaîne (on retrouve le pentane) ; les
positions 2 et 3 sont équivalentes par symétrie : un seul isomère, le \res{\text{2-méthylbutane}}.
\item Chaîne de 3 : deux méthyles, nécessairement sur le carbone central : le \res{\text{2,2-diméthylpropane}}.
\end{itemize}
Une chaîne de 3 ne peut pas porter d'éthyle (cela recréerait une chaîne de 4) : il n'existe donc que
\res{\text{trois squelettes}}.
\editionEtape{3. Les cinq squelettes de C6H14}
Même démarche : chaîne de 6, l'\res{\text{hexane}} ; chaîne de 5 avec un méthyle en 2 ou en 3 (positions non
équivalentes) : \res{\text{2-méthylpentane}} et \res{\text{3-méthylpentane}} ; chaîne de 4 avec deux méthyles : sur le
même carbone (\res{\text{2,2-diméthylbutane}}) ou sur deux carbones centraux (\res{\text{2,3-diméthylbutane}}).
Au total, \res{\text{cinq squelettes}}, dessinés page suivante.
\end{blocEdition}

\editionPage{defi-cor-11-2}{D11 · Corrigé (suite)}{Combien d'isomères ?}{\hyperref[defi-ex-11]{Énoncé p.~\pageref*{defi-ex-11}} · Démarche et vérifications}
\begin{blocEdition}{defi-cor-11-2}
\label{figedition-defi-cor-11-2-1}\schema[corrige]{defi-11-cor}
\editionEtape{4. Démasquer les doublons par le nom}
Placer un groupe éthyle en position 2 d'une chaîne de quatre carbones semble donner un nouveau squelette ; mais en
cherchant la chaîne la plus longue de ce dessin, on trouve cinq carbones : c'est le \res{\text{3-méthylpentane}},
déjà compté. La nomenclature impose de nommer chaque squelette à partir de sa chaîne la plus longue : deux dessins qui
reçoivent \res{\text{le même nom}} représentent la même molécule, quelle que soit leur apparence.
\editionEtape{5. L'explosion combinatoire}
De 3 isomères pour \ce{C5H12}, on passe à 75 pour \ce{C10H22} et à plus de \num{366000} pour \ce{C20H42} : il devient
impossible de les dessiner tous, et deux représentations d'une même molécule peuvent sembler très différentes. D'où
les deux outils du chimiste : la \res{\text{formule topologique}}, assez compacte pour représenter de grandes
molécules sans ambiguïté, et la \res{\text{nomenclature systématique}}, qui associe à chaque structure un nom unique
— et un seul.
\editionRetenir{Dénombrer des isomères exige une méthode : classer par longueur de chaîne principale décroissante,
placer les ramifications sans jamais allonger la chaîne, et nommer chaque candidat pour éliminer les doublons. Le
nombre d'isomères croît si vite avec la taille des molécules que ce travail systématique est au cœur de la chimie
organique.}
\end{blocEdition}

\editionPage{defi-cor-12-1}{D12 · Corrigé}{Jusqu'où pousser un rendement ?}{\hyperref[defi-ex-12]{Énoncé p.~\pageref*{defi-ex-12}} · Démarche et vérifications}
\begin{blocEdition}{defi-cor-12-1}
\editionEtape{1. Le palier des deux tiers}
La transformation est non totale : elle s'arrête lorsque $Q_\mathrm{r}=K$, pas lorsque les réactifs sont épuisés. En
mélange équimolaire (chapitre~7), $\frac{x^2}{(1-x)^2}=4$ donne $\frac{x}{1-x}=2$, soit $\res{x=\dfrac23}$ : prolonger
le chauffage n'y change rien, l'état d'équilibre est atteint.
\editionEtape{2. L'équation du taux final}
Avec 1~mol d'acide et $r$~mol d'alcool, l'état final contient $1-x$, $r-x$, $x$ et $x$ (en mol). Les volumes se
simplifient dans $Q_\mathrm{r}$ (chapitre~7) :
\[\frac{x^2}{(1-x)(r-x)}=4\ \Longleftrightarrow\ x^2=4\left(r-(1+r)x+x^2\right)\ \Longleftrightarrow\
\res{3x^2-4(1+r)\,x+4r=0}.\]
\editionEtape{3. La racine utile}
Le discriminant vaut $\Delta=16(1+r)^2-48r=16\left(r^2-r+1\right)>0$. La racine
$x=\frac{4(1+r)+\sqrt{\Delta}}{6}$ dépasse 1 : elle est à rejeter. Reste
\[x=\frac{4(1+r)-4\sqrt{r^2-r+1}}{6}=\res{\frac{2}{3}\left(1+r-\sqrt{r^2-r+1}\right)}.\]
Pour $r=1$ : $x=\frac23(2-1)=\res{2/3}$ ; pour $r=2$ : $x=\frac23\left(3-\sqrt3\right)=\res{0{,}845}$, les 85~\% de
l'exercice~10 du chapitre~7.
\editionEtape{4. L'excès nécessaire pour 95 \%}
Le programme essaie $r=1{,}0$ ; $1{,}5$ ; $2{,}0$\,\dots{} et s'arrête au premier excès qui atteint 95~\% : il affiche
$\res{r=5{,}5}$, pour un taux de 0,950. Il faut donc engager cinq fois et demie plus d'alcool que d'acide — dont
l'essentiel devra être séparé puis recyclé — pour gagner dix points par rapport à $r=2$ : chaque pour cent
supplémentaire coûte de plus en plus cher.
\end{blocEdition}

\editionPage{defi-cor-12-2}{D12 · Corrigé (suite)}{Jusqu'où pousser un rendement ?}{\hyperref[defi-ex-12]{Énoncé p.~\pageref*{defi-ex-12}} · Démarche et vérifications}
\begin{blocEdition}{defi-cor-12-2}
\label{figedition-defi-cor-12-2-1}\schema[corrige]{defi-12-cor}
\editionEtape{5. Excès, soutirage ou catalyseur}
La courbe montre un rendement qui tend vers 1 sans jamais l'atteindre : l'excès est une stratégie à rendements
décroissants. \res{\text{Éliminer l'eau}} au fur et à mesure (distillation, Dean-Stark, tamis moléculaire) est plus
efficace : $Q_\mathrm{r}$ est maintenu en dessous de $K$ et le système évolue dans le sens direct jusqu'à épuisement du
réactif limitant, sans excès à recycler. Un \res{\text{catalyseur}}, lui, accélère l'aller et le retour de la même
façon : il fait atteindre le palier plus vite, mais $K$ — donc le palier — est inchangé.
\editionRetenir{Le rendement d'une transformation équilibrée se pilote avec les outils du chapitre~7 : l'excès d'un
réactif déplace l'état final selon une loi que l'on peut calculer, le retrait d'un produit rend la transformation
quasi totale, et le catalyseur n'agit que sur la durée. Une stratégie de synthèse combine ces leviers en pesant leur
coût.}
\end{blocEdition}

\editionPage{defi-cor-13-1}{D13 · Corrigé}{Le caillou coincé dans un pneu}{\hyperref[defi-ex-13]{Énoncé p.~\pageref*{defi-ex-13}} · Démarche et vérifications}
\begin{blocEdition}{defi-cor-13-1}
\editionEtape{1. Le centre de la roue}
$y_\mathrm{C}=R$ est constant et $x_\mathrm{C}=R\omega t=V\,t$ : le centre de la roue a un mouvement
\res{\text{rectiligne uniforme}} à la vitesse $V$, comme le vélo.
\editionEtape{2. La vitesse du caillou}
En dérivant :
\[v_x=R\omega\left(1-\cos\omega t\right),\qquad v_y=R\omega\sin\omega t.\]
Au contact du sol ($\omega t=0$, $2\pi$\,\dots), $\cos\omega t=1$ et $\sin\omega t=0$ : $\res{v=0}$. Au sommet de la roue
($\omega t=\pi$), $v_x=2R\omega$ et $v_y=0$ : $\res{v=2V=\qty{16}{m.s^{-1}}}$. Le point de contact est momentanément
immobile — c'est la définition du roulement sans glissement — et le sommet de la roue va deux fois plus vite que le vélo.
\editionEtape{3. L'accélération}
En dérivant une seconde fois :
\[a_x=R\omega^2\sin\omega t,\qquad a_y=R\omega^2\cos\omega t,\]
ce sont exactement les expressions de \code{ax} et \code{ay} dans le programme.
\editionEtape{4. Une norme constante}
\[a=\sqrt{a_x^2+a_y^2}=R\omega^2\sqrt{\sin^2\omega t+\cos^2\omega t}=R\omega^2=\frac{V^2}{R}=\frac{8{,}0^2}{0{,}35}\approx\res{\qty{183}{m.s^{-2}}},\]
soit près de $19\,g$ : c'est pour cela que les cailloux finissent par être éjectés des pneus.
\end{blocEdition}

\editionPage{defi-cor-13-2}{D13 · Corrigé (suite)}{Le caillou coincé dans un pneu}{\hyperref[defi-ex-13]{Énoncé p.~\pageref*{defi-ex-13}} · Démarche et vérifications}
\begin{blocEdition}{defi-cor-13-2}
\label{figedition-defi-cor-13-2-1}\schema[corrige]{defi-13-cor}
\editionEtape{5. Toujours vers le centre}
Le vecteur $\vect{MC}$ a pour coordonnées $x_\mathrm{C}-x=R\sin\omega t$ et $y_\mathrm{C}-y=R\cos\omega t$. On a donc $\vec{a}=\omega^2\,\vect{MC}$ : l'accélération est \res{\text{toujours
dirigée vers le centre de la roue}}. Pourtant, dans le référentiel terrestre,
le mouvement n'est pas circulaire : la trajectoire est une cycloïde, et la vitesse varie de 0 à $2V$. Dans le
référentiel du vélo, en translation rectiligne uniforme par rapport au sol, le centre C est immobile et le caillou
décrit un cercle de rayon $R$ à la vitesse $R\omega=V$ : \res{\text{mouvement circulaire uniforme}}, d'accélération
centripète $V^2/R$. L'accélération est la même dans les deux référentiels, car ils sont en translation rectiligne
uniforme l'un par rapport à l'autre.
\editionRetenir{La trajectoire et la vitesse dépendent du référentiel ; l'accélération, elle, est la même dans deux
référentiels en translation rectiligne uniforme l'un par rapport à l'autre. Décrire un mouvement exige donc de préciser
le référentiel, puis de dériver soigneusement les coordonnées.}
\end{blocEdition}

\editionPage{defi-cor-14-1}{D14 · Corrigé}{De quelle hauteur lancer le wagon ?}{\hyperref[defi-ex-14]{Énoncé p.~\pageref*{defi-ex-14}} · Démarche et vérifications}
\begin{blocEdition}{defi-cor-14-1}
\editionEtape{1. La vitesse au sommet}
L'énergie mécanique se conserve entre le départ (vitesse nulle, hauteur $h$) et le sommet (hauteur $2R$) :
\[m\,g\,h=\frac12\,m\,v_\mathrm{S}^2+m\,g\,(2R)\qquad\text{d'où}\qquad\res{v_\mathrm{S}^2=2\,g\,(h-2R)}.\]
\editionEtape{2. La deuxième loi au sommet}
Dans le référentiel terrestre, galiléen, au sommet, $\vec{P}$ et $\vec{N}$ sont dirigés vers le centre du cercle, selon
$\un$. La projection de la deuxième loi sur $\un$ donne
\[m\,g+N=m\,\frac{v_\mathrm{S}^2}{R}\qquad\text{soit}\qquad\res{\frac{N}{m\,g}=\frac{v_\mathrm{S}^2}{g\,R}-1},\]
qui est exactement la ligne \code{n\_haut = v2\_haut / (g * R) - 1}.
\editionEtape{3. La hauteur minimale}
Le rail ne peut que pousser : le wagon reste en contact tant que $N\geq0$, c'est-à-dire $v_\mathrm{S}^2\geq g\,R$. Avec
la question 1 : $2\,g\,(h-2R)\geq g\,R$, soit
\[\res{h\geq\frac52\,R=\qty{15}{m}}.\]
Pour $h=\qty{12}{m}$, le programme affiche $N/(m\,g)=-1$ : une valeur négative est impossible. Le wagon n'atteindrait
même pas le sommet en contact avec le rail : il décrocherait avant et retomberait.
\editionEtape{4. Au point bas}
Au point bas, $v^2=2\,g\,h$ ; $\vec{N}$ est dirigé vers le haut (vers le centre) et $\vec{P}$ vers le bas. Selon $\un$ :
\[N-m\,g=m\,\frac{2\,g\,h}{R}\qquad\text{soit}\qquad\res{\frac{N}{m\,g}=1+\frac{2\,h}{R}}.\]
Pour $h=\frac52\,R$, on trouve $\res{N=6\,m\,g}$ : les passagers sont plaqués sur leur siège avec une force égale à six
fois leur poids.
\end{blocEdition}

\editionPage{defi-cor-14-2}{D14 · Corrigé (suite)}{De quelle hauteur lancer le wagon ?}{\hyperref[defi-ex-14]{Énoncé p.~\pageref*{defi-ex-14}} · Démarche et vérifications}
\begin{blocEdition}{defi-cor-14-2}
\label{figedition-defi-cor-14-2-1}\schema[corrige]{defi-14-cor}
\editionEtape{5. Pourquoi les loopings ne sont pas circulaires}
Avec un looping circulaire, le wagon doit partir d'au moins $\frac52R$ pour franchir le sommet, et subit alors $6\,g$
au point bas, au-delà de la limite de confort d'environ $5\,g$ ; une marge de sécurité ($h=\qty{18}{m}$) porterait même
cette valeur à $7\,g$. Les constructeurs dessinent donc des boucles en forme de goutte : un grand rayon de courbure en
bas, là où la vitesse est la plus grande, réduit $v^2/R$ ; un petit rayon au sommet permet de franchir celui-ci à
vitesse modérée, tout en gardant $N\geq0$.
\editionRetenir{La deuxième loi de Newton, écrite dans le repère de Frenet, relie la force exercée par le support à la
vitesse et au rayon de courbure. Associée à la conservation de l'énergie, elle donne des conditions de sécurité
concrètes : vitesse minimale au sommet, force maximale au point bas.}
\end{blocEdition}

\editionPage{defi-cor-15-1}{D15 · Corrigé}{Lancer à 45° ?}{\hyperref[defi-ex-15]{Énoncé p.~\pageref*{defi-ex-15}} · Démarche et vérifications}
\begin{blocEdition}{defi-cor-15-1}
\editionEtape{1. Depuis le sol}
Avec $h=0$ : $z=t\left(v_0\sin\alpha-\frac12\,g\,t\right)$ s'annule à $t=\dfrac{2\,v_0\sin\alpha}{g}$, d'où
\[D=v_0\cos\alpha\times\frac{2\,v_0\sin\alpha}{g}=\res{\frac{v_0^2\sin(2\alpha)}{g}}.\]
$\sin(2\alpha)$ vaut au plus 1, pour $2\alpha=90^\circ$ : la portée est maximale à $\res{\alpha=45^\circ}$ et vaut
$\dfrac{13{,}5^2}{9{,}81}=\res{\qty{18,6}{m}}$.
\editionEtape{2. Depuis la hauteur h}
Le poids touche le sol quand $z=-\frac12\,g\,t^2+v_z\,t+h=0$. Le discriminant vaut $v_z^2+2\,g\,h>0$ et la racine positive
est
\[t=\frac{v_z+\sqrt{v_z^2+2\,g\,h}}{g},\]
exactement la ligne du programme ; la longueur du jet vaut alors $v_x\,t$.
\editionEtape{3. Un optimum plus rasant}
Le programme affiche $\res{42^\circ}$ et \qty{20,57}{m}. Le jet est plus long que depuis le sol (\qty{18,6}{m}), car le poids
tombe de \qty{2,10}{m} supplémentaires en avançant. Cette chute finale avance d'autant plus que la vitesse horizontale
$v_0\cos\alpha$ est grande : elle favorise les angles plus faibles que $45^\circ$, ce qui déplace l'optimum vers le bas.
\end{blocEdition}

\editionPage{defi-cor-15-2}{D15 · Corrigé (suite)}{Lancer à 45° ?}{\hyperref[defi-ex-15]{Énoncé p.~\pageref*{defi-ex-15}} · Démarche et vérifications}
\begin{blocEdition}{defi-cor-15-2}
\label{figedition-defi-cor-15-2-1}\schema[corrige]{defi-15-cor}
\editionEtape{4. Pourquoi 37° en compétition ?}
À vitesse égale, on obtient \qty{20,57}{m} à $42^\circ$ et \qty{20,07}{m} à $35^\circ$ : l'écart n'est que de \qty{0,5}{m}, car le
maximum est très plat (voir la courbe). Or un athlète pousse plus fort lorsqu'il lance plus horizontalement : ses
muscles travaillent dans une position plus favorable. Gagner seulement 2~\% de vitesse augmente le jet d'environ
4~\%, puisque la longueur varie à peu près comme $v_0^2$ : cela compense largement les quelques décimètres perdus en
abaissant l'angle. L'angle optimal réel, autour de $37^\circ$, résulte de ce compromis entre la physique du projectile et
la physiologie du lanceur.
\editionRetenir{Une règle comme « 45° pour aller le plus loin » n'est valable que dans son modèle : départ au niveau du
sol, vitesse indépendante de l'angle, frottements négligés. Changer une seule hypothèse — ici la hauteur de lâcher,
puis la vitesse qui dépend de l'angle — suffit à modifier la conclusion.}
\end{blocEdition}

\editionPage{defi-cor-16-1}{D16 · Corrigé}{Peser le trou noir de la Galaxie}{\hyperref[defi-ex-16]{Énoncé p.~\pageref*{defi-ex-16}} · Démarche et vérifications}
\begin{blocEdition}{defi-cor-16-1}
\editionEtape{1. Une balance cosmique}
Pour l'étoile et pour la Terre, la troisième loi donne $\dfrac{T^2}{a^3}=\dfrac{4\pi^2}{GM}$ et $\dfrac{T_\oplus^2}{a_\oplus^3}=\dfrac{4\pi^2}{GM_\odot}$.
En divisant la seconde par la première :
\[\frac{M}{M_\odot}=\frac{a^3/a_\oplus^3}{T^2/T_\oplus^2}=\res{\frac{a^3}{T^2}}\quad(a\text{ en ua},\ T\text{ en années}),\]
puisque $a_\oplus=\qty{1}{ua}$ et $T_\oplus=\qty{1}{an}$ : c'est la ligne \code{M = a**3 / T**2}.
\editionEtape{2. Un seul objet central}
Les trois étoiles donnent \num{4.25}, \num{4.33} et \num{4.29} millions de masses solaires : à quelques pour cent près,
\res{\text{la même masse}}, environ $\res{\num{4.3e6}\,M_\odot}$. Les trois orbites sont gouvernées par un même objet, très massif
et invisible.
\editionEtape{3. Masse et rayon}
$M=\num{4.25e6}\times\num{1.99e30}=\res{\qty{8,5e36}{kg}}$, et
\[R_\mathrm{S}=\frac{2GM}{c^2}=\frac{2\times\num{6.67e-11}\times\num{8.47e36}}{(\num{3.00e8})^2}=\qty{1,25e10}{m}=\res{\qty{0,084}{ua}}.\]
Le trou noir tiendrait largement dans l'orbite de Mercure.
\editionEtape{4. Une approche très rapprochée}
Au plus près, S2 passe à $a(1-e)=1031\times0{,}114=\res{\qty{118}{ua}}$ de l'objet central, soit environ \res{1400} fois
$R_\mathrm{S}$. Une masse de quatre millions de soleils contenue dans une région aussi petite ne peut être ni une étoile ni
un amas d'étoiles : seul un trou noir est compatible avec ces observations.
\end{blocEdition}

\editionPage{defi-cor-16-2}{D16 · Corrigé (suite)}{Peser le trou noir de la Galaxie}{\hyperref[defi-ex-16]{Énoncé p.~\pageref*{defi-ex-16}} · Démarche et vérifications}
\begin{blocEdition}{defi-cor-16-2}
\label{figedition-defi-cor-16-2-1}\schema[corrige]{defi-16-cor}
\editionEtape{5. Des vitesses extrêmes}
Aux deux extrémités du grand axe, la vitesse est perpendiculaire au rayon : la loi des aires donne
$r_\text{min}\,v_\text{max}=r_\text{max}\,v_\text{min}$, d'où
\[\frac{v_\text{max}}{v_\text{min}}=\frac{a(1+e)}{a(1-e)}=\frac{1{,}886}{0{,}114}=\res{16{,}5}\qquad\text{et}\qquad
v_\text{min}=\frac{7650}{16{,}5}\approx\res{\qty{460}{km.s^{-1}}}.\]
Au plus près, S2 se déplace à 2,5~\% de la vitesse de la lumière : c'est l'une des étoiles les plus rapides connues.
\editionRetenir{La troisième loi de Kepler, établie pour les planètes, s'applique à toutes les orbites : mesurer la
période et la taille d'une orbite suffit à peser l'astre central, qu'il s'agisse du Soleil, de Jupiter ou d'un trou
noir. L'accord entre plusieurs orbites prouve qu'un même objet les gouverne.}
\end{blocEdition}

\editionPage{defi-cor-17-1}{D17 · Corrigé}{La clepsydre}{\hyperref[defi-ex-17]{Énoncé p.~\pageref*{defi-ex-17}} · Démarche et vérifications}
\begin{blocEdition}{defi-cor-17-1}
\editionEtape{1. Conservation du volume}
Pendant $\mathrm{d}t$, le vase perd le volume $\pi\,r(h)^2\,|\mathrm{d}h|$, qui est sorti par l'orifice à la vitesse $v$ :
$\pi\,r^2\,|\mathrm{d}h|=s\,\sqrt{2\,g\,h}\;\mathrm{d}t$. Le niveau baisse, $\mathrm{d}h<0$, d'où
\[u=-\frac{\mathrm{d}h}{\mathrm{d}t}=\res{\frac{s\,\sqrt{2\,g\,h}}{\pi\,r(h)^2}}.\]
\editionEtape{2. Le cylindre}
Avec $r=R$, posons $k=\dfrac{s}{\pi R^2}\sqrt{\dfrac{g}{2}}$. Si $\sqrt{h}=\sqrt{H}-k\,t$, alors
$\dfrac{1}{2\sqrt{h}}\,\dfrac{\mathrm{d}h}{\mathrm{d}t}=-k$, soit $\dfrac{\mathrm{d}h}{\mathrm{d}t}=-2k\sqrt{h}=-\dfrac{s\,\sqrt{2\,g\,h}}{\pi R^2}$ :
c'est bien la relation de la question~1, et $h(0)=H$. Le cylindre est vide quand $\sqrt{H}=k\,t$, soit
\[T_\mathrm{cyl}=\frac{\sqrt{H}}{k}=\res{\frac{\pi R^2}{s}\sqrt{\frac{2H}{g}}}.\]
La vitesse de descente, proportionnelle à $\sqrt{h}$, diminue au cours de la vidange : des graduations régulières
marqueraient des durées \res{\text{de plus en plus longues}}.
\editionEtape{3. Le profil idéal}
Si $u$ est constante, $\pi\,r^2\,u=s\,\sqrt{2\,g}\,\sqrt{h}$ : $r^2$ est proportionnel à $\sqrt{h}$, donc $r^4$ à $h$. Comme $r=R$
pour $h=H$ : $\res{h=H\,(r/R)^4}$, ou $r=R\,(h/H)^{1/4}$. Le vase s'évase vite près du fond, puis lentement : $r=\qty{8,41}{cm}$
à mi-hauteur et encore \qty{5,08}{cm} à \qty{2}{cm} du fond. Pour une vidange en 12~heures,
\[u=\frac{\qty{30}{cm}}{\qty{12}{h}}=\res{\qty{2,5}{cm.h^{-1}}}=\qty{6,94e-6}{m.s^{-1}}\quad\text{et, au sommet,}\quad
s=\frac{\pi R^2\,u}{\sqrt{2\,g\,H}}=\res{\qty{9,0e-8}{m^2}}.\]
L'orifice n'a qu'un diamètre $d=2\sqrt{s/\pi}=\res{\qty{0,34}{mm}}$ : un trou d'aiguille.
\editionEtape{4. Lecture du programme}
Dans la clepsydre, le niveau passe de 30 à 22,5 ; 15 ; 7,5 puis \qty{0}{cm} : il baisse de \res{\qty{7,5}{cm}} toutes les trois
heures, soit \qty{2,5}{cm.h^{-1}}, et le vase est vide au bout de 12~heures. Dans le cylindre, les baisses successives
valent 7,0 ; 6,1 ; 5,2 puis \qty{4,2}{cm} : elles \res{\text{diminuent}}, et il reste encore $\qty{7,5}{cm}=H/4$ au bout de
12~heures, conformément à la question~2.
\end{blocEdition}

\editionPage{defi-cor-17-2}{D17 · Corrigé (suite)}{La clepsydre}{\hyperref[defi-ex-17]{Énoncé p.~\pageref*{defi-ex-17}} · Démarche et vérifications}
\begin{blocEdition}{defi-cor-17-2}
\label{figedition-defi-cor-17-2-1}\schema[corrige]{defi-17-cor}
\editionEtape{5. Deux fois plus lent}
En remplaçant $s$ par $\dfrac{\pi R^2\,u}{\sqrt{2\,g\,H}}$ dans $T_\mathrm{cyl}$ :
\[T_\mathrm{cyl}=\frac{\sqrt{2\,g\,H}}{u}\sqrt{\frac{2H}{g}}=\res{\frac{2H}{u}}=\qty{24}{h},\]
le double des 12~heures de la clepsydre. Les deux vases ont le même rayon au sommet : le niveau y descend d'abord à la
même vitesse, puis celui du cylindre ralentit. Un tronc de cône évasé vers le haut, facile à tailler dans la pierre,
réduit la section vers le fond et compense en partie ce ralentissement : c'est une approximation du profil idéal.
\editionRetenir{Deux lois simples, la conservation du débit et la formule de Torricelli, suffisent à prévoir la vidange
d'un récipient de forme quelconque. Inverser le problème, en choisissant la forme pour imposer le mouvement voulu, est
une démarche d'ingénieur.}
\end{blocEdition}

\editionPage{defi-cor-18-1}{D18 · Corrigé}{Faire le vide}{\hyperref[defi-ex-18]{Énoncé p.~\pageref*{defi-ex-18}} · Démarche et vérifications}
\begin{blocEdition}{defi-cor-18-1}
\editionEtape{1. Le premier coup de pompe}
Pendant l'aspiration, l'air du récipient garde la même quantité de matière et la même température, mais occupe le volume
$V+v$ : la loi de Mariotte donne $p_1\,(V+v)=p_0\,V$, d'où
\[p_1=p_0\,\frac{V}{V+v}=1013\times\frac{1{,}0}{1{,}050}=\res{\qty{965}{hPa}}.\]
Le refoulement de l'air du cylindre ne modifie pas l'air resté dans le récipient, isolé par la soupape.
\editionEtape{2. Une suite géométrique}
Le même raisonnement s'applique à chaque coup : $p_{k+1}=q\,p_k$ avec $q=\dfrac{V}{V+v}=\dfrac{1}{1{,}05}$. La suite est
géométrique de raison $q$ et de premier terme $p_0$ :
\[\res{p_k=p_0\left(\frac{V}{V+v}\right)^{k}}.\]
Chaque coup retire la même \textbf{fraction} de l'air restant, environ 4,8~\% : la pression diminue de moins en moins.
\editionEtape{3. Combien de coups ?}
On cherche le plus petit entier $k$ tel que $p_0\,q^k<\qty{1}{hPa}$, soit $k\ln q<\ln\dfrac{1}{1013}$. Comme $\ln q<0$, le sens de
l'inégalité change :
\[k>\frac{\ln1013}{\ln1{,}05}=141{,}8\qquad\text{d'où}\qquad \res{k=142\ \text{coups}},\]
comme l'affiche le programme, avec $p_{142}=\qty{0,992}{hPa}$. Atteindre \qty{20}{hPa} ne demande que 81 coups : les derniers
hectopascals sont les plus longs à gagner.
\editionEtape{4. Gonfler n'est pas l'inverse de vider}
Chaque coup fait entrer dans le pneu la même quantité d'air, celle de \qty{0,10}{L} sous \qty{1,0}{bar} : la pression augmente
de $\dfrac{0{,}10}{1{,}5}\times\qty{1,0}{bar}$ à chaque coup, $p_k=\res{p_0\left(1+k\,\dfrac{v}{V}\right)}$. La suite est
\textbf{arithmétique} : 30 coups suffisent pour atteindre \qty{3,0}{bar}. Pour le vide, chaque coup retire une
\textbf{fraction} de l'air restant : la suite est géométrique.
\end{blocEdition}

\editionPage{defi-cor-18-2}{D18 · Corrigé (suite)}{Faire le vide}{\hyperref[defi-ex-18]{Énoncé p.~\pageref*{defi-ex-18}} · Démarche et vérifications}
\begin{blocEdition}{defi-cor-18-2}
\label{figedition-defi-cor-18-2-1}\schema[corrige]{defi-18-cor}
\editionEtape{5. Le volume mort impose une limite}
À l'aspiration, l'air du récipient ($p_k$, $V$) et celui du volume mort ($p_0$, $v_\mathrm{m}$) occupent ensemble $V+v+v_\mathrm{m}$.
La quantité totale se conserve : $p_{k+1}\,(V+v+v_\mathrm{m})=p_k\,V+p_0\,v_\mathrm{m}$, c'est la ligne du programme. Une pression
qui ne change plus vérifie $p\,(v+v_\mathrm{m})=p_0\,v_\mathrm{m}$, d'où
\[p_\mathrm{lim}=p_0\,\frac{v_\mathrm{m}}{v+v_\mathrm{m}}=1013\times\frac{0{,}50}{50{,}5}=\res{\qty{10,0}{hPa}}.\]
La pression s'en approche sans jamais descendre en dessous : la condition \code{p > 1} reste vraie et la boucle s'arrête à
\code{k = 1000}. Le programme affiche \res{\text{1000 10.03}}. Les pompes à vide réelles utilisent plusieurs étages pour
contourner cette limite.
\editionRetenir{À température constante, une suite de détentes identiques divise la pression par un même facteur :
la pression décroît exponentiellement avec le nombre de coups. Un détail technique, le volume mort, suffit à fixer une
limite que la modélisation permet de prévoir.}
\end{blocEdition}

\editionPage{defi-cor-19-1}{D19 · Corrigé}{L'expérience de Joule}{\hyperref[defi-ex-19]{Énoncé p.~\pageref*{defi-ex-19}} · Démarche et vérifications}
\begin{blocEdition}{defi-cor-19-1}
\editionEtape{1. Un travail qui échauffe}
Le système reçoit un \res{\text{travail}} : les palettes, entraînées par les masses, exercent sur l'eau des forces dont
les points d'application se déplacent. Aucun corps plus chaud ne le touche et l'on néglige les échanges avec l'air :
$Q\approx0$. Le premier principe donne $\Delta U=W>0$ : l'énergie interne augmente, ici sous forme d'agitation des
molécules d'eau, et la température s'élève.
\editionEtape{2. Le travail reçu}
Pour une descente des deux masses ($m=\qty{26,0}{kg}$), l'énergie mécanique perdue vaut
\[m\,g\,h-\tfrac12\,m\,v^2=26{,}0\times9{,}81\times1{,}60-\tfrac12\times26{,}0\times0{,}060^2=\qty{408,1}{J}-\qty{0,047}{J}.\]
Le terme cinétique ne représente qu'un dix-millième de l'énergie potentielle : il est \res{\text{négligeable}}. Pour
20 descentes : $W=\res{\qty{8,16e3}{J}}$.
\editionEtape{3. L'élévation prévue}
$\Delta U=W$ avec $\Delta U=(m_\mathrm{eau}\,c_\mathrm{eau}+C)\,\Delta T$ :
\[\Delta T=\frac{W}{m_\mathrm{eau}\,c_\mathrm{eau}+C}=\frac{\num{8.16e3}}{6{,}00\times4180+400}=\res{\qty{0,320}{K}}.\]
Une élévation d'un tiers de degré seulement : il faut un thermomètre très précis.
\editionEtape{4. L'équivalent mécanique de la calorie}
La masse d'eau de même capacité thermique que le récipient vaut $\dfrac{C}{c_\mathrm{eau}}=\dfrac{400}{4180}=\qty{0,0957}{kg}$, soit
\res{\qty{95,7}{g}}. Le système reçoit donc, en calories, $(6000+95{,}7)\times0{,}313=\qty{1908}{cal}$, en échange d'un travail de
\qty{8,16e3}{J} :
\[\frac{\num{8.16e3}\ \text{J}}{\num{1908}\ \text{cal}}=\res{\qty{4,28}{J}\ \text{par calorie}},\]
soit \qty{2,3}{\percent} de plus que la valeur actuelle.
\end{blocEdition}

\editionPage{defi-cor-19-2}{D19 · Corrigé (suite)}{L'expérience de Joule}{\hyperref[defi-ex-19]{Énoncé p.~\pageref*{defi-ex-19}} · Démarche et vérifications}
\begin{blocEdition}{defi-cor-19-2}
\label{figedition-defi-cor-19-2-1}\schema[corrige]{defi-19-cor}
\editionEtape{5. Les pertes, la répétition et le premier principe}
L'élévation mesurée correspond à une énergie interne de $25\,480\times0{,}313=\qty{7,98e3}{J}$ : environ \qty{0,19}{kJ}, soit
2~\%, n'a pas servi à échauffer l'eau. Une partie du travail est perdue dans les \res{\text{frottements}} des poulies et des
axes, et l'eau cède de l'énergie à l'air pendant la longue durée de l'expérience. L'élévation mesurée est donc trop
faible, ce qui \res{\text{surestime}} le nombre de joules par calorie ; Joule corrigeait ces pertes et obtenait
\qty{4,16}{J}.

Une seule descente n'élève la température que de \qty{0,016}{K}, en deçà de la précision des thermomètres : répéter vingt
fois la chute \res{\text{multiplie l'effet}} mesuré par vingt.

L'expérience montre qu'un travail et un transfert thermique produisent le même effet sur un système et qu'ils
s'échangent dans un rapport fixe : ce sont deux modes de transfert d'une même grandeur, l'énergie. C'est le fondement
du premier principe, $\Delta U=W+Q$, et la raison pour laquelle on mesure aujourd'hui les transferts thermiques en
joules.
\editionRetenir{Le premier principe n'est pas une définition : il repose sur des mesures, comme celles de Joule, qui
établissent l'équivalence entre travail et transfert thermique. Une mesure précise d'un très petit effet demande
d'amplifier ce dernier et d'évaluer les pertes.}
\end{blocEdition}

\editionPage{defi-cor-20-1}{D20 · Corrigé}{Café : le lait avant ou après ?}{\hyperref[defi-ex-20]{Énoncé p.~\pageref*{defi-ex-20}} · Démarche et vérifications}
\begin{blocEdition}{defi-cor-20-1}
\editionEtape{1. La température du mélange}
Pendant le mélange, le système \{café, lait\} n'échange ni travail ni transfert thermique avec l'extérieur :
$\Delta U=0$. Avec la même capacité thermique massique $c$ :
\[m_\mathrm{c}\,c\,(T-T_\mathrm{c})+m_\mathrm{l}\,c\,(T-T_\mathrm{l})=0\qquad\text{d'où}\qquad \res{T=\frac{m_\mathrm{c}T_\mathrm{c}+m_\mathrm{l}T_\mathrm{l}}{m_\mathrm{c}+m_\mathrm{l}}},\]
moyenne des températures pondérée par les masses : c'est la fonction \code{melanger}.
\editionEtape{2. Option A : le lait d'abord}
La fonction \code{refroidir} est la solution de l'équation de Newton, $T(t)=T_\mathrm{e}+(T_0-T_\mathrm{e})\,\mathrm{e}^{-t/\tau}$. Le
mélange part de $\dfrac{200\times90+50\times5}{250}=\qty{73}{\degreeCelsius}$ ; au bout de \qty{600}{s}, avec
$\mathrm{e}^{-600/1500}=0{,}670$ :
\[T_\mathrm{A}=20+(73-20)\times0{,}670=\res{\qty{55,53}{\degreeCelsius}}.\]
\editionEtape{3. Option B : le lait à la fin}
Le café seul atteint $20+70\times0{,}670=\qty{66,92}{\degreeCelsius}$, puis le lait l'amène à
\[T_\mathrm{B}=\frac{200\times66{,}92+50\times5}{250}=\res{\qty{54,54}{\degreeCelsius}}.\]
Verser le lait \res{\text{à la fin}} donne un café plus frais d'environ \qty{1}{\degreeCelsius}. D'après la loi de Newton, le flux
perdu est proportionnel à l'écart avec la pièce : pendant l'attente, le café seul, plus chaud que le mélange, perd son
énergie plus vite. Inversement, pour garder un café chaud, mieux vaut ajouter le lait tout de suite.
\end{blocEdition}

\editionPage{defi-cor-20-2}{D20 · Corrigé (suite)}{Café : le lait avant ou après ?}{\hyperref[defi-ex-20]{Énoncé p.~\pageref*{defi-ex-20}} · Démarche et vérifications}
\begin{blocEdition}{defi-cor-20-2}
\label{figedition-defi-cor-20-2-1}\schema[corrige]{defi-20-cor}
\editionEtape{4. Une formule générale}
Notons $M=m_\mathrm{c}+m_\mathrm{l}$ et $k=\mathrm{e}^{-t/\tau}$. En développant
$T_\mathrm{A}=T_\mathrm{e}+\left(\dfrac{m_\mathrm{c}T_\mathrm{c}+m_\mathrm{l}T_\mathrm{l}}{M}-T_\mathrm{e}\right)k$ et
$T_\mathrm{B}=\dfrac{m_\mathrm{c}\left(T_\mathrm{e}+(T_\mathrm{c}-T_\mathrm{e})\,k\right)+m_\mathrm{l}T_\mathrm{l}}{M}$, les termes en $T_\mathrm{c}$ se compensent :
\[T_\mathrm{A}-T_\mathrm{B}=\res{\frac{m_\mathrm{l}}{M}\,(T_\mathrm{e}-T_\mathrm{l})\,(1-k)}=0{,}20\times15\times0{,}330=\qty{0,99}{K}.\]
L'écart est positif dès que le lait est plus froid que la pièce. Si le lait est à la température de la pièce, les deux
options \res{\text{se valent}}. L'écart maximal, pour une attente très longue, vaut $0{,}20\times15=\res{\qty{3,0}{K}}$.
\editionEtape{5. Discuter l'hypothèse}
Le temps caractéristique $\tau=\dfrac{C}{h\,S}$ est proportionnel à la capacité thermique. Avec le lait, le contenu de la tasse
refroidit plus lentement : le mélange A garde davantage sa chaleur, et le café seul de l'option B refroidit plus vite.
Les deux effets \res{\text{renforcent}} la conclusion, en partie compensés par l'augmentation de la surface mouillée.
\editionRetenir{Avec la loi de Newton, les pertes sont proportionnelles à l'écart de température : l'ordre des opérations
compte. Une modélisation simple permet de trancher une question de la vie courante et d'en prévoir les limites.}
\end{blocEdition}

\editionPage{defi-cor-21-1}{D21 · Corrigé}{Pourquoi entend-on les basses du voisin ?}{\hyperref[defi-ex-21]{Énoncé p.~\pageref*{defi-ex-21}} · Démarche et vérifications}
\begin{blocEdition}{defi-cor-21-1}
\editionEtape{1. Six décibels par octave}
En doublant la fréquence : $20\log(m_\mathrm{s}\times2f)-47=20\log(m_\mathrm{s}\,f)-47+20\log2$, soit
\[R(2f)=R(f)+\res{\qty{6,0}{dB}}.\]
La sortie le confirme : 42,2 ; 48,2 ; 54,2 ; 60,2 ; 66,3 ; 72,3~dB, soit environ \qty{6}{dB} d'écart entre deux lignes.
\editionEtape{2. Une conversation de l'autre côté}
Le niveau transmis vaut $L-R$. Derrière le béton : $60-54{,}2=\res{\qty{5,8}{dB}}$, inaudible. Derrière la cloison, dont
l'indice vaut $R=20\log(20\times500)-47=\qty{33}{dB}$ : $60-33=\res{\qty{27}{dB}}$, comme un murmure, bien audible dans une
pièce calme. La masse de la paroi fait toute la différence.
\editionEtape{3. Les basses passent, les aigus non}
À \qty{63}{Hz}, le béton n'atténue que $R=20\log(230\times63)-47=\qty{36,2}{dB}$ : il reste $90-36{,}2=\res{\qty{53,8}{dB}}$. À
\qty{4000}{Hz}, $R=\qty{72,3}{dB}$ et il ne reste que \res{\qty{7,7}{dB}}. De l'autre côté du mur, on n'entend plus que les
\res{\text{basses}} : les parois atténuent d'autant moins un son qu'il est grave.
\editionEtape{4. Alourdir le mur ?}
Gagner \qty{20}{dB} demande $20\log k=20$, soit une masse multipliée par \res{10} : un mur de béton d'un mètre
d'épaisseur. Ce n'est pas réaliste ; on utilise plutôt des doubles parois séparées par un isolant fibreux, qui
dépassent largement la loi de masse.
\end{blocEdition}

\editionPage{defi-cor-21-2}{D21 · Corrigé (suite)}{Pourquoi entend-on les basses du voisin ?}{\hyperref[defi-ex-21]{Énoncé p.~\pageref*{defi-ex-21}} · Démarche et vérifications}
\begin{blocEdition}{defi-cor-21-2}
\label{figedition-defi-cor-21-2-1}\schema[corrige]{defi-21-cor}
\editionEtape{5. Le maillon faible}
Mur : $\tau_\mathrm{m}=10^{-5{,}42}=\num{3.8e-6}$ ; porte : $R=20\log(15\times500)-47=\qty{30,5}{dB}$ et $\tau_\mathrm{p}=\num{8.9e-4}$. En
moyenne, sur \qty{10}{m^2} :
\[\tau=\frac{8\times\num{3.8e-6}+2\times\num{8.9e-4}}{10}=\num{1.81e-4}\qquad\text{d'où}\qquad R=-10\log\tau=\res{\qty{37,4}{dB}}.\]
La porte, qui n'occupe qu'un cinquième de la surface, transmet 98~\% de l'énergie : l'isolation de l'ensemble chute
de 54 à \qty{37}{dB}. En acoustique, une paroi isole comme son point le plus faible : porte, fenêtre ou simple fente.
\editionRetenir{Les décibels traduisent des rapports d'intensité : une paroi qui atténue de \qty{54}{dB} ne laisse passer
que quatre millionièmes de l'intensité, mais la moindre faiblesse ruine ce résultat. La loi de masse explique
pourquoi les basses traversent les murs.}
\end{blocEdition}

\editionPage{defi-cor-22-1}{D22 · Corrigé}{Des franges dans la tache de diffraction}{\hyperref[defi-ex-22]{Énoncé p.~\pageref*{defi-ex-22}} · Démarche et vérifications}
\begin{blocEdition}{defi-cor-22-1}
\editionEtape{1. La tache de diffraction d'une fente}
\[L=\frac{2\lambda D}{a}=\frac{2\times\num{633e-9}\times2{,}0}{\num{0.10e-3}}=\res{\qty{25,3}{mm}}.\]
\editionEtape{2. L'interfrange}
$i=\dfrac{\lambda D}{b}=\res{\qty{2,53}{mm}}$, et $\dfrac{L}{i}=\dfrac{2\lambda D/a}{\lambda D/b}=\res{\dfrac{2b}{a}}=\dfrac{2\times0{,}50}{0{,}10}=10$ : la tache
centrale s'étend sur dix interfranges.
\editionEtape{3. Une frange manquante}
Le terme de diffraction s'annule pour $\sin u=0$ avec $u\neq0$ : le premier zéro correspond à $u=\pi$, soit
$x=\dfrac{\lambda D}{a}=\qty{12,66}{mm}$, le bord de la tache centrale. La frange brillante d'ordre $k$ est en $x_k=k\,i$ ; pour
$k=5$, $x_5=5\times\qty{2,532}{mm}=\qty{12,66}{mm}$ : elle tombe exactement sur le zéro de diffraction et \res{\text{disparaît}}.
Il reste les franges d'ordres $-4$ à 4, soit \res{\text{9 franges brillantes}} dans la tache centrale.
\editionEtape{4. Lecture du programme}
Le programme affiche, pour chaque frange $k$, sa position en millimètres et son intensité relative : 1 ; 0,875 ; 0,573 ;
0,255 ; 0,055 et 0. Au centre de chaque frange brillante, le terme d'interférences vaut 1 : l'intensité est celle de la
\res{\text{diffraction}}, qui décroît du centre vers le bord de la tache. L'intensité nulle de la frange 5 confirme la
question~3.
\end{blocEdition}

\editionPage{defi-cor-22-2}{D22 · Corrigé (suite)}{Des franges dans la tache de diffraction}{\hyperref[defi-ex-22]{Énoncé p.~\pageref*{defi-ex-22}} · Démarche et vérifications}
\begin{blocEdition}{defi-cor-22-2}
\label{figedition-defi-cor-22-2-1}\schema[corrige]{defi-22-cor}
\editionEtape{5. Des fentes plus fines}
Avec $a=\qty{0,05}{mm}$, $\dfrac{L}{i}=\dfrac{2b}{a}=20$ : la tache centrale s'élargit à \qty{50,6}{mm} et la frange d'ordre 10 tombe sur
le premier zéro. On compte les franges d'ordres $-9$ à 9, soit \res{\text{19 franges brillantes}}. L'interfrange, lui, ne change
pas. Mais chaque fente, deux fois plus étroite, laisse passer deux fois moins de lumière, répartie sur une tache deux
fois plus large : la figure est \res{\text{beaucoup plus sombre}}, ce qui gêne l'observation et la mesure.
\editionRetenir{Diffraction et interférences coexistent toujours : la figure observée derrière des fentes d'Young est
une figure d'interférences modulée par la figure de diffraction de chaque fente. Comparer les deux échelles, $\lambda D/a$ et
$\lambda D/b$, suffit à prévoir son aspect.}
\end{blocEdition}

\editionPage{defi-cor-23-1}{D23 · Corrigé}{Le passage d'une voiture}{\hyperref[defi-ex-23]{Énoncé p.~\pageref*{defi-ex-23}} · Démarche et vérifications}
\begin{blocEdition}{defi-cor-23-1}
\editionEtape{1. Loin de l'auditeur}
$\cos\theta=\dfrac{-x}{\sqrt{x^2+b^2}}$ tend vers 1 quand $x\to-\infty$ (la voiture arrive droit sur l'auditeur, vue de loin) et vers
$-1$ quand $x\to+\infty$. On retrouve les relations à une dimension :
\[f_\mathrm{app}=900\times\frac{340}{315}=\res{\qty{971,4}{Hz}}\qquad\text{et}\qquad f_\text{éloi}=900\times\frac{340}{365}=\res{\qty{838,4}{Hz}}.\]
Le programme s'en approche : 971,0~Hz pour $t=\qty{-4}{s}$ et 838,6~Hz pour $t=\qty{4}{s}$.
\editionEtape{2. Au plus près}
En $x=0$, la vitesse de la voiture est perpendiculaire à la direction de l'auditeur : $\cos\theta=0$ et
$f_\mathrm{R}=\res{f_\mathrm{E}=\qty{900}{Hz}}$. Ce son parcourt $b=\qty{10}{m}$ : il est reçu à $t=\dfrac{10}{340}=\res{\qty{0,03}{s}}$.
\editionEtape{3. Le temps de propagation}
Le son émis à la date $t$ parcourt la distance $d$ avant d'atteindre l'auditeur : il est reçu à la date $t+\dfrac{d}{v}$,
d'autant plus tard que la voiture est loin. Le passage entendu est donc légèrement \res{\text{retardé}} par rapport au passage
réel.
\editionEtape{4. Une dissymétrie}
$f_\mathrm{app}-f_\mathrm{E}=\qty{71,4}{Hz}$ mais $f_\mathrm{E}-f_\text{éloi}=\qty{61,6}{Hz}$ : la hausse dépasse la baisse. La moyenne des deux
valeurs, \qty{904,9}{Hz}, \res{\text{n'est pas}} $f_\mathrm{E}$. Il faut éliminer $f_\mathrm{E}$ entre les deux relations, comme dans le
programme de l'essentiel : $v_\mathrm{E}=v\,\dfrac{f_\mathrm{app}-f_\text{éloi}}{f_\mathrm{app}+f_\text{éloi}}=\res{\qty{25,0}{m.s^{-1}}}$, puis
$f_\mathrm{E}=f_\mathrm{app}\,\dfrac{v-v_\mathrm{E}}{v}=\qty{900}{Hz}$.
\end{blocEdition}

\editionPage{defi-cor-23-2}{D23 · Corrigé (suite)}{Le passage d'une voiture}{\hyperref[defi-ex-23]{Énoncé p.~\pageref*{defi-ex-23}} · Démarche et vérifications}
\begin{blocEdition}{defi-cor-23-2}
\label{figedition-defi-cor-23-2-1}\schema[corrige]{defi-23-cor}
\editionEtape{5. Une transition plus ou moins brusque}
L'angle $\theta$ change notablement lorsque la voiture passe de $x\approx-b$ à $x\approx+b$ : $\cos\theta$ va alors de 0,71 à $-0,71$.
Cette traversée dure environ $\dfrac{2b}{v_\mathrm{E}}=\res{\qty{0,8}{s}}$ pour $b=\qty{10}{m}$, mais \res{\qty{4}{s}} pour $b=\qty{50}{m}$ : loin de la
route, la fréquence glisse progressivement ; tout près, elle chute presque d'un coup, ce qui donne le « vroum »
caractéristique d'une voiture qui passe à côté de nous.
\editionRetenir{Les relations à une dimension décrivent les phases d'approche et d'éloignement lointaines. Pour en
déduire une vitesse, il faut utiliser les deux fréquences extrêmes et éliminer la fréquence émise : le décalage
Doppler n'est pas symétrique.}
\end{blocEdition}

\editionPage{defi-cor-24-1}{D24 · Corrigé}{Où placer l'œil ? Le cercle oculaire}{\hyperref[defi-ex-24]{Énoncé p.~\pageref*{defi-ex-24}} · Démarche et vérifications}
\begin{blocEdition}{defi-cor-24-1}
\editionEtape{1. Position du cercle oculaire}
Pour l'oculaire, l'objectif est un objet réel situé en $\overline{OA}=-(f'_1+f'_2)$. La relation de conjugaison donne
\[\frac{1}{\overline{OA'}}=\frac{1}{f'_2}-\frac{1}{f'_1+f'_2}=\frac{f'_1}{f'_2\,(f'_1+f'_2)}\qquad\text{d'où}\qquad \res{\overline{OA'}=x=\frac{f'_2\,(f'_1+f'_2)}{f'_1}}>0 :\]
l'image est réelle, derrière l'oculaire, un peu au-delà de son foyer image. Pour $f'_2=\qty{10}{mm}$ :
$x=\dfrac{10\times710}{700}=\res{\qty{10,1}{mm}}$.
\editionEtape{2. Diamètre du cercle oculaire}
$\gamma=\dfrac{\overline{OA'}}{\overline{OA}}=\dfrac{f'_2\,(f'_1+f'_2)/f'_1}{-(f'_1+f'_2)}=-\dfrac{f'_2}{f'_1}$ : l'image de l'objectif est renversée et réduite dans le
rapport $f'_2/f'_1$, d'où $\res{d'=D\,\dfrac{f'_2}{f'_1}=\dfrac{D}{G}}$.
\editionEtape{3. Lecture du programme}
Les colonnes donnent la distance focale de l'oculaire, le grossissement, la distance $x$ et le diamètre $d'$. Plus la focale
de l'oculaire est courte, plus le grossissement est grand et plus le cercle oculaire est \res{\text{petit et proche}} de
l'oculaire : \qty{4,0}{mm} à \qty{42}{mm} de distance avec l'oculaire de \qty{40}{mm}, mais \qty{0,4}{mm} à \qty{4}{mm} seulement avec celui
de \qty{4}{mm}. L'œil doit alors presque toucher le verre pour placer sa pupille dans le cercle oculaire.
\end{blocEdition}

\editionPage{defi-cor-24-2}{D24 · Corrigé (suite)}{Où placer l'œil ? Le cercle oculaire}{\hyperref[defi-ex-24]{Énoncé p.~\pageref*{defi-ex-24}} · Démarche et vérifications}
\begin{blocEdition}{defi-cor-24-2}
\label{figedition-defi-cor-24-2-1}\schema[corrige]{defi-24-cor}
\editionEtape{4. Ne pas perdre de lumière}
Toute la lumière entre dans l'œil si $d'\leqslant\qty{7}{mm}$, soit $\dfrac{D}{G}\leqslant7$ : $\res{G\geqslant10}$. L'oculaire de plus grande focale
utilisable a donc $f'_2=\dfrac{700}{10}=\res{\qty{70}{mm}}$ ; au-delà, une partie du faisceau tombe à côté de la pupille.
\editionEtape{5. Éclat et grossissement maximal}
Avec l'oculaire de \qty{10}{mm}, $d'=\qty{1,0}{mm}$ contre \qty{4,0}{mm} avec celui de \qty{40}{mm} : l'image est $\left(\dfrac{1{,}0}{4{,}0}\right)^2$,
soit \res{\text{16 fois moins lumineuse}}, car la même lumière est étalée sur une image 4 fois plus grande. Pour $G=2D$,
$d'=\dfrac{D}{2D}=\res{\qty{0,5}{mm}}$ : en deçà, l'image devient trop sombre et la diffraction par la petite pupille de sortie la rend
floue. La règle du grossissement maximal utile revient à garder un cercle oculaire d'au moins un demi-millimètre.
\editionRetenir{Le cercle oculaire relie le grossissement à la lumière reçue par l'œil : $d'=D/G$. Un grand objectif permet
à la fois un fort grossissement et une image lumineuse ; c'est pourquoi le diamètre est la caractéristique principale
d'un instrument d'astronomie.}
\end{blocEdition}

\editionPage{defi-cor-25-1}{D25 · Corrigé}{Mesurer la constante de Planck avec des LED}{\hyperref[defi-ex-25]{Énoncé p.~\pageref*{defi-ex-25}} · Démarche et vérifications}
\begin{blocEdition}{defi-cor-25-1}
\editionEtape{1. Une droite passant par l'origine}
L'énergie reçue par un électron se retrouve dans le photon émis : $e\,U_\mathrm{s}=\dfrac{h\,c}{\lambda}$, soit
\[U_\mathrm{s}=\res{\frac{h\,c}{e}\times\frac{1}{\lambda}}.\]
La tension de seuil est proportionnelle à $1/\lambda$ : le graphique de $U_\mathrm{s}$ en fonction de $1/\lambda$ est une droite passant
par l'origine, de coefficient directeur $\dfrac{h\,c}{e}$.
\editionEtape{2. Le programme}
La liste \code{x} contient les valeurs de $1/\lambda$ en \unit{m^{-1}} : le facteur \code{1e9} convertit les nanomètres. Les
variables \code{num} et \code{den} réalisent une \res{\text{régression linéaire}} par la méthode des moindres carrés : \code{a} est le
coefficient directeur de la droite qui passe au plus près des cinq points, en \unit{V.m}. Le programme l'affiche
multiplié par $10^6$, c'est-à-dire en \unit{V.\micro m} : $a=\qty{1,203}{V.\micro m}=\qty{1,203e-6}{V.m}$.
\editionEtape{3. La constante de Planck}
$h=\dfrac{a\,e}{c}=\dfrac{\num{1.203e-6}\times\num{1.602e-19}}{\num{2.998e8}}=\res{\qty{6,43e-34}{J.s}}$, soit un écart relatif de
$\dfrac{6{,}43-6{,}626}{6{,}626}\approx\res{\qty{-3,0}{\percent}}$ : un résultat remarquable pour une expérience aussi simple.
\end{blocEdition}

\editionPage{defi-cor-25-2}{D25 · Corrigé (suite)}{Mesurer la constante de Planck avec des LED}{\hyperref[defi-ex-25]{Énoncé p.~\pageref*{defi-ex-25}} · Démarche et vérifications}
\begin{blocEdition}{defi-cor-25-2}
\label{figedition-defi-cor-25-2-1}\schema[corrige]{defi-25-cor}
\editionEtape{4. Un écart systématique}
Avec $\dfrac{h\,c}{e}=\qty{1239,8}{V.nm}$, on attend 1,32 ; 1,97 ; 2,10 ; 2,36 et \qty{2,64}{V}. Toutes les mesures sont inférieures
d'environ 3~\% : l'écart est \res{\text{systématique}}. La « tension de seuil » est difficile à définir, le courant augmentant
progressivement ; une LED émet sur une bande de longueurs d'onde, pas une seule ; enfin, l'agitation thermique fournit aux
électrons une énergie de l'ordre de $kT\approx\qty{0,026}{eV}$ qui facilite le passage du courant. Une erreur systématique
affecte toutes les mesures dans le même sens : la régression ne l'élimine pas.
\editionEtape{5. L'intérêt de l'infrarouge}
La LED infrarouge élargit l'intervalle des valeurs de $1/\lambda$ (de 1,06 à \qty{2,13}{\micro m^{-1}}) : le coefficient directeur est
\res{\text{mieux déterminé}}. Sa lumière, invisible à l'œil, est détectée par le capteur d'un smartphone, qui la montre sur
l'écran sous forme d'une tache violacée.
\editionRetenir{Une constante fondamentale peut se mesurer avec un matériel modeste : il suffit d'un modèle (ici, un
électron donne un photon), d'une relation linéaire et de plusieurs mesures exploitées par régression. L'analyse critique
des écarts distingue erreurs aléatoires et erreurs systématiques.}
\end{blocEdition}

\editionPage{defi-cor-26-1}{D26 · Corrigé}{Où passe la moitié de l'énergie ?}{\hyperref[defi-ex-26]{Énoncé p.~\pageref*{defi-ex-26}} · Démarche et vérifications}
\begin{blocEdition}{defi-cor-26-1}
\editionEtape{1. L'énergie fournie par le générateur}
En fin de charge, $u_C=E$ et $q=C\,E=\num{1.0e-3}\times10=\res{\qty{0,010}{C}}$. Toute cette charge a traversé le générateur sous la
tension constante $E$ : il a fourni le travail électrique $W_\mathrm{g}=E\,q=\res{C\,E^2}=\qty{0,10}{J}$.
\editionEtape{2. Un rendement de 50~\%}
Le condensateur stocke $\frac12\,C\,E^2=\res{\qty{0,050}{J}}$ : la \res{\text{moitié}} de l'énergie fournie a été perdue, et le rendement
de la charge vaut \res{\qty{50}{\percent}}, quelle que soit la valeur des composants.
\editionEtape{3. L'énergie dissipée par effet Joule}
Avec $i=\dfrac{E}{R}\,\mathrm{e}^{-t/\tau}$, la puissance dissipée vaut $R\,i^2=\dfrac{E^2}{R}\,\mathrm{e}^{-2t/\tau}$. Une primitive de $\mathrm{e}^{-2t/\tau}$ est
$-\dfrac{\tau}{2}\,\mathrm{e}^{-2t/\tau}$, d'où
\[W_\mathrm{J}=\frac{E^2}{R}\left[-\frac{\tau}{2}\,\mathrm{e}^{-2t/\tau}\right]_0^{+\infty}=\frac{E^2}{R}\times\frac{RC}{2}=\res{\frac12\,C\,E^2}.\]
La résistance s'élimine : l'énergie perdue ne dépend pas de $R$. On vérifie le bilan : $W_\mathrm{g}=\frac12CE^2+\frac12CE^2$.
\editionEtape{4. Lecture du programme}
Pour \qty{10}{\ohm}, \qty{100}{\ohm} et \qty{1000}{\ohm}, la somme vaut \qty{0,05}{J}, égale à l'énergie stockée. À $t=10\tau$, la puissance
dissipée n'est plus que $\mathrm{e}^{-20}\approx\num{2e-9}$ fois sa valeur initiale : la suite de la somme est \res{\text{négligeable}}.
\end{blocEdition}

\editionPage{defi-cor-26-2}{D26 · Corrigé (suite)}{Où passe la moitié de l'énergie ?}{\hyperref[defi-ex-26]{Énoncé p.~\pageref*{defi-ex-26}} · Démarche et vérifications}
\begin{blocEdition}{defi-cor-26-2}
\label{figedition-defi-cor-26-2-1}\schema[corrige]{defi-26-cor}
\editionEtape{5. Charger progressivement}
Dans chaque étape, la tension aux bornes de la résistance part de $E/2$ et décroît exponentiellement : le calcul de la
question~3 s'applique avec $E/2$ au lieu de $E$. Chaque étape dissipe $\frac12\,C\left(\frac{E}{2}\right)^2=\qty{0,0125}{J}$, soit
\qty{0,025}{J} au total, contre \qty{0,050}{J} en une seule fois. Le rendement devient $\dfrac{0{,}050}{0{,}050+0{,}025}=\res{\qty{66,7}{\percent}}$.
Avec $N$ étapes, chacune dissipe $\frac12\,C\left(\frac{E}{N}\right)^2$ : les pertes totales valent $\dfrac{C\,E^2}{2N}$ et tendent vers zéro quand
$N$ augmente, par exemple 90,9~\% de rendement avec 10 étapes. Les chargeurs de supercondensateurs augmentent ainsi la
tension progressivement.
\editionRetenir{Brancher brutalement un condensateur sur un générateur de tension gaspille toujours la moitié de l'énergie,
quelle que soit la résistance : une faible résistance dissipe plus fort, mais moins longtemps. Les pertes viennent de
la différence de tension entre le générateur et le condensateur ; la réduire réduit les pertes.}
\end{blocEdition}

\editionPage{defi-cor-27-1}{D27 · Corrigé}{Mesurer plus pour mesurer mieux ?}{\hyperref[defi-ex-27]{Énoncé p.~\pageref*{defi-ex-27}} · Démarche et vérifications}
\begin{blocEdition}{defi-cor-27-1}
\editionEtape{1. Pourquoi diviser par la racine de N}
La moyenne s'écrit $\bar m=\dfrac{1}{N}\,S$, avec $S=m_1+m_2+\dots+m_N$. Les mesures étant indépendantes, la règle des sommes
donne $u(S)^2=s^2+s^2+\dots+s^2=N\,s^2$, soit $u(S)=\sqrt N\,s$. Avec la constante $k=\dfrac1N$ :
\[u(\bar m)=\frac{\sqrt N\,s}{N}=\res{\frac{s}{\sqrt N}}.\]
\editionEtape{2. Lecture du programme}
Le programme simule \num{2000} séries de 40 mesures tirées autour de \qty{0,452}{s} avec l'écart-type $s$, et range la moyenne
de chaque série dans la liste \code{moyennes}. L'écart-type de ces \num{2000} moyennes, \qty{0,00205}{s}, est bien égal à
$\dfrac{s}{\sqrt{40}}=\qty{0,00206}{s}$ : les moyennes fluctuent \res{\sqrt{40}\approx6{,}3} fois moins que les mesures elles-mêmes.
\editionEtape{3. Un retard finit toujours par se voir}
Sans retard, $t=\sqrt{\dfrac{2\times1{,}000}{9{,}81}}=\res{\qty{0,4515}{s}}$. Chaque mesure vaut $t+\delta$ plus une fluctuation aléatoire : la
moyenne tend vers $t+\delta$, avec l'incertitude-type $s/\sqrt N$. L'écart normalisé avec la durée réelle vaut donc
\[z=\frac{\delta}{s/\sqrt N}=\frac{0{,}010\,\sqrt N}{0{,}013}.\]
Le retard est détecté dès que $z>2$, soit $\sqrt N>2{,}6$ : $N>6{,}76$, donc \res{\text{dès 7 mesures}}. Avec 40 mesures,
$z=\res{4{,}9}$ : la moyenne est nettement incompatible avec la durée réelle.
\editionEtape{4. Plus précis, mais pas plus juste}
Avec $\bar t=0{,}4515+0{,}010=\qty{0,4615}{s}$, $g=\dfrac{2\times1{,}000}{0{,}4615^2}=\res{\qty{9,39}{m.s^{-2}}}$, avec
$\dfrac{u(g)}{g}=\sqrt{0{,}0020^2+\left(\dfrac{2\times0{,}0020}{0{,}4615}\right)^2}=0{,}89~\%$, soit $u(g)=\res{\qty{0,084}{m.s^{-2}}}$, et $z=\dfrac{0{,}42}{0{,}084}=\res{5{,}0}$.
Multiplier les mesures a rendu le résultat plus \textbf{précis}, mais pas plus \textbf{juste} : l'erreur systématique
reste entière. En revanche, la faible incertitude la rend \res{\text{visible}}.
\end{blocEdition}

\editionPage{defi-cor-27-2}{D27 · Corrigé (suite)}{Mesurer plus pour mesurer mieux ?}{\hyperref[defi-ex-27]{Énoncé p.~\pageref*{defi-ex-27}} · Démarche et vérifications}
\begin{blocEdition}{defi-cor-27-2}
\editionEtape{5. Séparer le retard et la pesanteur}
$t=\sqrt{\dfrac{2h}{g}}=\sqrt{\dfrac2g}\,\sqrt h$, d'où $t_\mathrm{mes}=\sqrt{\dfrac2g}\,\sqrt h+\delta$ : une \res{\text{fonction affine}} de $\sqrt h$. Avec
$\sqrt{0{,}20}=0{,}447$, le coefficient directeur vaut $a=\dfrac{0{,}4615-0{,}2119}{1-0{,}447}=\qty{0,4515}{s.m^{-1/2}}$ (la régression linéaire
sur les cinq points donne la même valeur). L'ordonnée à l'origine vaut $0{,}4615-0{,}4515=\res{\qty{0,010}{s}}=\delta$, et
$g=\dfrac{2}{a^2}=\res{\qty{9,81}{m.s^{-2}}}$.
\label{figedition-defi-cor-27-2-1}\schema[corrige]{defi-27-cor}
\editionRetenir{La moyenne de $N$ mesures réduit les erreurs aléatoires comme $1/\sqrt N$, mais ne réduit pas une erreur
systématique : plus l'incertitude diminue, plus un biais devient visible. Faire varier un paramètre, ici la hauteur,
et exploiter une droite permet de séparer le biais de la grandeur cherchée.}
\end{blocEdition}
